\documentclass[UTF8,fontset=fandol,openany]{ctexbook}
% ctexbook defaults are retained; only page geometry and technical support are added.
\usepackage[a4paper,inner=18mm,outer=16mm,top=18mm,bottom=19mm,headheight=15pt]{geometry}
\usepackage{amsmath,amssymb,amsthm,mathtools}
\usepackage{graphicx,tikz,adjustbox}
\usetikzlibrary{arrows.meta,positioning,calc,fit,shapes.geometric,decorations.pathreplacing,decorations.pathmorphing}
\usepackage{array,longtable,booktabs,multirow,tabularx}
\usepackage[normalem]{ulem}
\usepackage{xeCJKfntef}
\xeCJKDeclareCharClass{CJK}{"2460 -> "2473}
\usepackage{enumitem}
\usepackage{hyperref}
\hypersetup{unicode=true,hidelinks,bookmarksnumbered=true,pdfcreator={XeLaTeX / ctexbook}}
\usepackage{bookmark}
\setcounter{tocdepth}{1}
\setcounter{secnumdepth}{2}
\let\cleardoublepage\clearpage
% Large nonbreaking TikZ figures may move to the next page. Keep ordinary
% paragraph spacing instead of stretching the preceding page vertically.
\raggedbottom
\setlength{\emergencystretch}{3em}
\newcommand{\SourceTag}[1]{\textbf{【#1】}}
\newcommand{\SourceFootnote}[1]{\par\smallskip\noindent{\footnotesize #1}\par\smallskip}
\newcommand{\InsertDiagram}[1]{\par\smallskip\begin{center}\begin{adjustbox}{max width=\linewidth,max totalheight=.78\textheight}\input{#1}\end{adjustbox}\end{center}\smallskip}
\newcommand{\OriginalPanel}[1]{\par\smallskip\noindent\textbf{#1}\par}
\newcommand{\BookEditionDate}{2026 年 11 月 20 日}
\newcommand{\BookCover}[4]{%
  \begin{titlepage}\centering
  \vspace*{20mm}
  {\zihao{2}\bfseries #1\par}
  \vspace{10mm}
  {\zihao{-0}\bfseries 刑法\par}
  \vspace{8mm}
  {\zihao{2}#2\par}
  \vspace{12mm}
  {\zihao{4}#3\par}
  \vfill
  {\zihao{-4}#4\par}
  \vspace{9mm}
  {\zihao{-4}\BookEditionDate\par}
  \vspace*{14mm}
  \end{titlepage}}

\begin{document}
\setbox0=\vbox{\hsize=4.340cm\parindent=0pt\linespread{1}\fontsize{10}{13.50}\selectfont\bfseries\centering\noindent 行为时\par}
\typeout{TCELL:p025_1:0:\the\ht0:\the\dp0}
\setbox0=\vbox{\hsize=4.340cm\parindent=0pt\linespread{1}\fontsize{10}{13.50}\selectfont\bfseries\centering\noindent 审判时\par}
\typeout{TCELL:p025_1:1:\the\ht0:\the\dp0}
\setbox0=\vbox{\hsize=6.540cm\parindent=0pt\linespread{1}\fontsize{10}{13.50}\selectfont\bfseries\centering\noindent 处理结论\par}
\typeout{TCELL:p025_1:2:\the\ht0:\the\dp0}
\setbox0=\vbox{\hsize=4.340cm\parindent=0pt\linespread{1}\fontsize{10}{13.50}\selectfont\raggedright\noindent 无司法解释\par}
\typeout{TCELL:p025_1:3:\the\ht0:\the\dp0}
\setbox0=\vbox{\hsize=4.340cm\parindent=0pt\linespread{1}\fontsize{10}{13.50}\selectfont\raggedright\noindent 有司法解释\par}
\typeout{TCELL:p025_1:4:\the\ht0:\the\dp0}
\setbox0=\vbox{\hsize=6.540cm\parindent=0pt\linespread{1}\fontsize{10}{13.50}\selectfont\raggedright\noindent 按照司法解释处理②\par}
\typeout{TCELL:p025_1:5:\the\ht0:\the\dp0}
\setbox0=\vbox{\hsize=4.340cm\parindent=0pt\linespread{1}\fontsize{10}{13.50}\selectfont\raggedright\noindent 有司法解释\par}
\typeout{TCELL:p025_1:6:\the\ht0:\the\dp0}
\setbox0=\vbox{\hsize=4.340cm\parindent=0pt\linespread{1}\fontsize{10}{13.50}\selectfont\raggedright\noindent 新司法解释\par}
\typeout{TCELL:p025_1:7:\the\ht0:\the\dp0}
\setbox0=\vbox{\hsize=6.540cm\parindent=0pt\linespread{1}\fontsize{10}{13.50}\selectfont\raggedright\noindent 按照“从旧兼从轻”原则处理\par}
\typeout{TCELL:p025_1:8:\the\ht0:\the\dp0}
\setbox0=\vbox{\hsize=3.740cm\parindent=0pt\linespread{1}\fontsize{10}{13.50}\selectfont\bfseries\centering\noindent 身份与罪名\par}
\typeout{TCELL:p037_1:0:\the\ht0:\the\dp0}
\setbox0=\vbox{\hsize=3.440cm\parindent=0pt\linespread{1}\fontsize{10}{13.50}\selectfont\bfseries\centering\noindent 贪钱\par}
\typeout{TCELL:p037_1:1:\the\ht0:\the\dp0}
\setbox0=\vbox{\hsize=4.540cm\parindent=0pt\linespread{1}\fontsize{10}{13.50}\selectfont\bfseries\centering\noindent 收钱\par}
\typeout{TCELL:p037_1:2:\the\ht0:\the\dp0}
\setbox0=\vbox{\hsize=3.240cm\parindent=0pt\linespread{1}\fontsize{10}{13.50}\selectfont\bfseries\centering\noindent 挪钱\par}
\typeout{TCELL:p037_1:3:\the\ht0:\the\dp0}
\setbox0=\vbox{\hsize=3.740cm\parindent=0pt\linespread{1}\fontsize{10}{13.50}\selectfont\raggedright\noindent 国家工作人员\par}
\typeout{TCELL:p037_1:4:\the\ht0:\the\dp0}
\setbox0=\vbox{\hsize=3.440cm\parindent=0pt\linespread{1}\fontsize{10}{13.50}\selectfont\raggedright\noindent 贪污罪\par}
\typeout{TCELL:p037_1:5:\the\ht0:\the\dp0}
\setbox0=\vbox{\hsize=4.540cm\parindent=0pt\linespread{1}\fontsize{10}{13.50}\selectfont\raggedright\noindent 受贿罪\par}
\typeout{TCELL:p037_1:6:\the\ht0:\the\dp0}
\setbox0=\vbox{\hsize=3.240cm\parindent=0pt\linespread{1}\fontsize{10}{13.50}\selectfont\raggedright\noindent 挪用公款罪\par}
\typeout{TCELL:p037_1:7:\the\ht0:\the\dp0}
\setbox0=\vbox{\hsize=3.740cm\parindent=0pt\linespread{1}\fontsize{10}{13.50}\selectfont\raggedright\noindent 非国家工作人员\par}
\typeout{TCELL:p037_1:8:\the\ht0:\the\dp0}
\setbox0=\vbox{\hsize=3.440cm\parindent=0pt\linespread{1}\fontsize{10}{13.50}\selectfont\raggedright\noindent 职务侵占罪\par}
\typeout{TCELL:p037_1:9:\the\ht0:\the\dp0}
\setbox0=\vbox{\hsize=4.540cm\parindent=0pt\linespread{1}\fontsize{10}{13.50}\selectfont\raggedright\noindent 非国家工作人员受贿罪\par}
\typeout{TCELL:p037_1:10:\the\ht0:\the\dp0}
\setbox0=\vbox{\hsize=3.240cm\parindent=0pt\linespread{1}\fontsize{10}{13.50}\selectfont\raggedright\noindent 挪用资金罪\par}
\typeout{TCELL:p037_1:11:\the\ht0:\the\dp0}
\setbox0=\vbox{\hsize=4.040cm\parindent=0pt\linespread{1}\fontsize{10}{13.50}\selectfont\bfseries\centering\noindent 事项\par}
\typeout{TCELL:p075_2:0:\the\ht0:\the\dp0}
\setbox0=\vbox{\hsize=5.540cm\parindent=0pt\linespread{1}\fontsize{10}{13.50}\selectfont\bfseries\centering\noindent 故意犯罪\par}
\typeout{TCELL:p075_2:1:\the\ht0:\the\dp0}
\setbox0=\vbox{\hsize=5.640cm\parindent=0pt\linespread{1}\fontsize{10}{13.50}\selectfont\bfseries\centering\noindent 过失犯罪\par}
\typeout{TCELL:p075_2:2:\the\ht0:\the\dp0}
\setbox0=\vbox{\hsize=4.040cm\parindent=0pt\linespread{1}\fontsize{10}{13.50}\selectfont\raggedright\noindent 成立条件\par}
\typeout{TCELL:p075_2:3:\the\ht0:\the\dp0}
\setbox0=\vbox{\hsize=5.540cm\parindent=0pt\linespread{1}\fontsize{10}{13.50}\selectfont\raggedright\noindent 一般不要求发生实害结果\par}
\typeout{TCELL:p075_2:4:\the\ht0:\the\dp0}
\setbox0=\vbox{\hsize=5.640cm\parindent=0pt\linespread{1}\fontsize{10}{13.50}\selectfont\raggedright\noindent 要求发生实害结果\par}
\typeout{TCELL:p075_2:5:\the\ht0:\the\dp0}
\setbox0=\vbox{\hsize=4.040cm\parindent=0pt\linespread{1}\fontsize{10}{13.50}\selectfont\raggedright\noindent 主客观相一致原则\par}
\typeout{TCELL:p075_2:6:\the\ht0:\the\dp0}
\setbox0=\vbox{\hsize=5.540cm\parindent=0pt\linespread{1}\fontsize{10}{13.50}\selectfont\raggedright\noindent 要求遵守主客观相一致原则\par}
\typeout{TCELL:p075_2:7:\the\ht0:\the\dp0}
\setbox0=\vbox{\hsize=5.640cm\parindent=0pt\linespread{1}\fontsize{10}{13.50}\selectfont\raggedright\noindent 不存在遵守该原则的必要\par}
\typeout{TCELL:p075_2:8:\the\ht0:\the\dp0}
\setbox0=\vbox{\hsize=4.040cm\parindent=0pt\linespread{1}\fontsize{10}{13.50}\selectfont\raggedright\noindent 主观恶性程度\par}
\typeout{TCELL:p075_2:9:\the\ht0:\the\dp0}
\setbox0=\vbox{\hsize=5.540cm\parindent=0pt\linespread{1}\fontsize{10}{13.50}\selectfont\raggedright\noindent 较重，因此刑罚也较重\par}
\typeout{TCELL:p075_2:10:\the\ht0:\the\dp0}
\setbox0=\vbox{\hsize=5.640cm\parindent=0pt\linespread{1}\fontsize{10}{13.50}\selectfont\raggedright\noindent 较轻，因此刑罚也较轻\par}
\typeout{TCELL:p075_2:11:\the\ht0:\the\dp0}
\setbox0=\vbox{\hsize=4.040cm\parindent=0pt\linespread{1}\fontsize{10}{13.50}\selectfont\raggedright\noindent 未完成形态\par}
\typeout{TCELL:p075_2:12:\the\ht0:\the\dp0}
\setbox0=\vbox{\hsize=5.540cm\parindent=0pt\linespread{1}\fontsize{10}{13.50}\selectfont\raggedright\noindent 有（包括间接故意）\par}
\typeout{TCELL:p075_2:13:\the\ht0:\the\dp0}
\setbox0=\vbox{\hsize=5.640cm\parindent=0pt\linespread{1}\fontsize{10}{13.50}\selectfont\raggedright\noindent 无，只有成立与否问题\par}
\typeout{TCELL:p075_2:14:\the\ht0:\the\dp0}
\setbox0=\vbox{\hsize=4.040cm\parindent=0pt\linespread{1}\fontsize{10}{13.50}\selectfont\raggedright\noindent 共同犯罪\par}
\typeout{TCELL:p075_2:15:\the\ht0:\the\dp0}
\setbox0=\vbox{\hsize=5.540cm\parindent=0pt\linespread{1}\fontsize{10}{13.50}\selectfont\raggedright\noindent 有\par}
\typeout{TCELL:p075_2:16:\the\ht0:\the\dp0}
\setbox0=\vbox{\hsize=5.640cm\parindent=0pt\linespread{1}\fontsize{10}{13.50}\selectfont\raggedright\noindent 无\par}
\typeout{TCELL:p075_2:17:\the\ht0:\the\dp0}
\setbox0=\vbox{\hsize=4.040cm\parindent=0pt\linespread{1}\fontsize{10}{13.50}\selectfont\raggedright\noindent 处罚\par}
\typeout{TCELL:p075_2:18:\the\ht0:\the\dp0}
\setbox0=\vbox{\hsize=5.540cm\parindent=0pt\linespread{1}\fontsize{10}{13.50}\selectfont\raggedright\noindent 处罚故意犯罪是原则\par}
\typeout{TCELL:p075_2:19:\the\ht0:\the\dp0}
\setbox0=\vbox{\hsize=5.640cm\parindent=0pt\linespread{1}\fontsize{10}{13.50}\selectfont\raggedright\noindent 处罚过失犯罪是例外\par}
\typeout{TCELL:p075_2:20:\the\ht0:\the\dp0}
\setbox0=\vbox{\hsize=3.140cm\parindent=0pt\linespread{1}\fontsize{10}{13.50}\selectfont\bfseries\centering\noindent 关于犯罪既遂的条件\par}
\typeout{TCELL:p081_1:0:\the\ht0:\the\dp0}
\setbox0=\vbox{\hsize=3.340cm\parindent=0pt\linespread{1}\fontsize{10}{13.50}\selectfont\bfseries\centering\noindent 因果关系\par}
\typeout{TCELL:p081_1:1:\the\ht0:\the\dp0}
\setbox0=\vbox{\hsize=8.740cm\parindent=0pt\linespread{1}\fontsize{10}{13.50}\selectfont\bfseries\centering\noindent 范例\par}
\typeout{TCELL:p081_1:2:\the\ht0:\the\dp0}
\setbox0=\vbox{\hsize=3.140cm\parindent=0pt\linespread{1}\fontsize{10}{13.50}\selectfont\centering\noindent 客观要件\par}
\typeout{TCELL:p081_1:3:\the\ht0:\the\dp0}
\setbox0=\vbox{\hsize=3.340cm\parindent=0pt\linespread{1}\fontsize{10}{13.50}\selectfont\raggedright\noindent 没有因果关系，不构成故意犯罪既遂\par}
\typeout{TCELL:p081_1:4:\the\ht0:\the\dp0}
\setbox0=\vbox{\hsize=8.740cm\parindent=0pt\linespread{1}\fontsize{10}{13.50}\selectfont\raggedright\noindent （人参果案）孙悟空想毒杀猪八戒，给其有毒的人参果。猪八戒吃时卡在喉咙，噎死。毒药没有进入胃里起作用，毒杀行为与死亡没有因果关系，不构成故意杀人罪既遂，而是未遂①（2023 年试题）\par}
\typeout{TCELL:p081_1:5:\the\ht0:\the\dp0}
\setbox0=\vbox{\hsize=3.340cm\parindent=0pt\linespread{1}\fontsize{10}{13.50}\selectfont\raggedright\noindent 具有因果关系，进入主观要件$\to$\par}
\typeout{TCELL:p081_1:6:\the\ht0:\the\dp0}
\setbox0=\vbox{\hsize=3.140cm\parindent=0pt\linespread{1}\fontsize{10}{13.50}\selectfont\centering\noindent 主观要件\par}
\typeout{TCELL:p081_1:7:\the\ht0:\the\dp0}
\setbox0=\vbox{\hsize=3.340cm\parindent=0pt\linespread{1}\fontsize{10}{13.50}\selectfont\raggedright\noindent 以为是 A 死法，实际是 B 死法。\\{}[结论]构成故意犯罪既遂\par}
\typeout{TCELL:p081_1:8:\the\ht0:\the\dp0}
\setbox0=\vbox{\hsize=8.740cm\parindent=0pt\linespread{1}\fontsize{10}{13.50}\selectfont\raggedright\noindent （推下井案）甲将乙推入井里，欲淹死乙，井里没水，摔死了乙。客观上，甲的行为与乙的死亡存在因果关系。主观上，甲关于具体因果流程的认识错误不重要。甲构成故意杀人罪既遂。（2011 年第 53 题）\\{}[总结]人参果案与推下井案的相同点是，都以为是 A 死法导致死亡，实际是 B 死法导致死亡。区别在于，前者在客观要件中没有因果关系；后者在客观要件中有因果关系\par}
\typeout{TCELL:p081_1:9:\the\ht0:\the\dp0}
\setbox0=\vbox{\hsize=3.140cm\parindent=0pt\linespread{1}\fontsize{10}{13.50}\selectfont\bfseries\centering\noindent 防卫过当的成立条件\par}
\typeout{TCELL:p098_1:0:\the\ht0:\the\dp0}
\setbox0=\vbox{\hsize=3.540cm\parindent=0pt\linespread{1}\fontsize{10}{13.50}\selectfont\bfseries\centering\noindent 具体要求\par}
\typeout{TCELL:p098_1:1:\the\ht0:\the\dp0}
\setbox0=\vbox{\hsize=8.540cm\parindent=0pt\linespread{1}\fontsize{10}{13.50}\selectfont\bfseries\centering\noindent 范例\par}
\typeout{TCELL:p098_1:2:\the\ht0:\the\dp0}
\setbox0=\vbox{\hsize=3.140cm\parindent=0pt\linespread{1}\fontsize{10}{13.50}\selectfont\raggedright\noindent 客观要件\par}
\typeout{TCELL:p098_1:3:\the\ht0:\the\dp0}
\setbox0=\vbox{\hsize=3.540cm\parindent=0pt\linespread{1}\fontsize{10}{13.50}\selectfont\raggedright\noindent 发生过当结果\par}
\typeout{TCELL:p098_1:4:\the\ht0:\the\dp0}
\setbox0=\vbox{\hsize=8.540cm\parindent=0pt\linespread{1}\fontsize{10}{13.50}\selectfont\raggedright\noindent 甲轻伤害乙，乙开枪反击，虽然手段过当，但只造成甲轻伤。对此不能认定为防卫过当。也即，即使手段过当，但没有过当结果，不构成防卫过当\par}
\typeout{TCELL:p098_1:5:\the\ht0:\the\dp0}
\setbox0=\vbox{\hsize=3.140cm\parindent=0pt\linespread{1}\fontsize{10}{13.50}\selectfont\raggedright\noindent 主观要件\par}
\typeout{TCELL:p098_1:6:\the\ht0:\the\dp0}
\setbox0=\vbox{\hsize=3.540cm\parindent=0pt\linespread{1}\fontsize{10}{13.50}\selectfont\raggedright\noindent 对过当结果至少有过失\par}
\typeout{TCELL:p098_1:7:\the\ht0:\the\dp0}
\setbox0=\vbox{\hsize=8.540cm\parindent=0pt\linespread{1}\fontsize{10}{13.50}\selectfont\raggedright\noindent 甲轻伤害乙，乙用木棍反击，竟致甲死亡。（1）此时，乙属于客观意义上的防卫过当。（2）到了主观要件，如果对甲的死亡有过失，则成立最终意义上的防卫过当，以过失致人死亡罪论处；如果没过失，属于意外事件，则不负刑事责任①\par}
\typeout{TCELL:p098_1:8:\the\ht0:\the\dp0}
\setbox0=\vbox{\hsize=1.840cm\parindent=0pt\linespread{1}\fontsize{10}{13.50}\selectfont\bfseries\centering\noindent 教唆犯的\\成立条件\par}
\typeout{TCELL:p146_1:0:\the\ht0:\the\dp0}
\setbox0=\vbox{\hsize=2.140cm\parindent=0pt\linespread{1}\fontsize{10}{13.50}\selectfont\centering\noindent 客观要件\par}
\typeout{TCELL:p146_1:1:\the\ht0:\the\dp0}
\setbox0=\vbox{\hsize=3.940cm\parindent=0pt\linespread{1}\fontsize{10}{13.50}\selectfont\centering\noindent 引起正犯制造违法事实\par}
\typeout{TCELL:p146_1:2:\the\ht0:\the\dp0}
\setbox0=\vbox{\hsize=7.040cm\parindent=0pt\linespread{1}\fontsize{10}{13.50}\selectfont\raggedright\noindent 引起正犯制造违法行为，是教唆犯的成立条件\par}
\typeout{TCELL:p146_1:3:\the\ht0:\the\dp0}
\setbox0=\vbox{\hsize=7.040cm\parindent=0pt\linespread{1}\fontsize{10}{13.50}\selectfont\raggedright\noindent 引起正犯制造违法结果，是教唆犯的既遂条件\par}
\typeout{TCELL:p146_1:4:\the\ht0:\the\dp0}
\setbox0=\vbox{\hsize=2.140cm\parindent=0pt\linespread{1}\fontsize{10}{13.50}\selectfont\centering\noindent 主观要件\par}
\typeout{TCELL:p146_1:5:\the\ht0:\the\dp0}
\setbox0=\vbox{\hsize=11.240cm\parindent=0pt\linespread{1}\fontsize{10}{13.50}\selectfont\centering\noindent 故意引起\par}
\typeout{TCELL:p146_1:6:\the\ht0:\the\dp0}
\setbox0=\vbox{\hsize=3.840cm\parindent=0pt\linespread{1}\fontsize{10}{13.50}\selectfont\bfseries\centering\noindent 罪名关系的类型\par}
\typeout{TCELL:p173_1:0:\the\ht0:\the\dp0}
\setbox0=\vbox{\hsize=5.940cm\parindent=0pt\linespread{1}\fontsize{10}{13.50}\selectfont\bfseries\centering\noindent 特点\par}
\typeout{TCELL:p173_1:1:\the\ht0:\the\dp0}
\setbox0=\vbox{\hsize=5.440cm\parindent=0pt\linespread{1}\fontsize{10}{13.50}\selectfont\bfseries\centering\noindent 范例\par}
\typeout{TCELL:p173_1:2:\the\ht0:\the\dp0}
\setbox0=\vbox{\hsize=3.840cm\parindent=0pt\linespread{1}\fontsize{10}{13.50}\selectfont\raggedright\noindent A 与 $-A$ 的对立排斥关系（男生与女生的逻辑关系）\par}
\typeout{TCELL:p173_1:3:\the\ht0:\the\dp0}
\setbox0=\vbox{\hsize=5.940cm\parindent=0pt\linespread{1}\fontsize{10}{13.50}\selectfont\raggedright\noindent 一个行为针对一个对象，不可能同时触犯这两个罪名，也即既不可能属于法条竞合，也不可能形成想象竞合\par}
\typeout{TCELL:p173_1:4:\the\ht0:\the\dp0}
\setbox0=\vbox{\hsize=5.440cm\parindent=0pt\linespread{1}\fontsize{10}{13.50}\selectfont\raggedright\noindent 甲欺骗乙回头看看，甲趁乙回头，拿走乙的财物。甲对乙，要么定诈骗罪，要么定盗窃罪；应定盗窃罪\par}
\typeout{TCELL:p173_1:5:\the\ht0:\the\dp0}
\setbox0=\vbox{\hsize=3.840cm\parindent=0pt\linespread{1}\fontsize{10}{13.50}\selectfont\raggedright\noindent A 与 A+B 的包容评价关系（本科生与硕士生的逻辑关系）\par}
\typeout{TCELL:p173_1:6:\the\ht0:\the\dp0}
\setbox0=\vbox{\hsize=5.940cm\parindent=0pt\linespread{1}\fontsize{10}{13.50}\selectfont\raggedright\noindent A+B 可以包容评价为 A。触犯 A+B，必然触犯 A。所以这种关系也是法条竞合关系，也被称为位阶关系\par}
\typeout{TCELL:p173_1:7:\the\ht0:\the\dp0}
\setbox0=\vbox{\hsize=5.440cm\parindent=0pt\linespread{1}\fontsize{10}{13.50}\selectfont\raggedright\noindent 甲想盗窃普通财物，实际盗窃到枪支。盗窃枪支罪可以包容评价为盗窃罪。甲构成盗窃罪既遂\par}
\typeout{TCELL:p173_1:8:\the\ht0:\the\dp0}
\setbox0=\vbox{\hsize=3.840cm\parindent=0pt\linespread{1}\fontsize{10}{13.50}\selectfont\raggedright\noindent A 与 B 的中立关系（男生与硕士生的逻辑关系）\par}
\typeout{TCELL:p173_1:9:\the\ht0:\the\dp0}
\setbox0=\vbox{\hsize=5.940cm\parindent=0pt\linespread{1}\fontsize{10}{13.50}\selectfont\raggedright\noindent A 与 B 不可能属于法条竞合关系，但有可能形成想象竞合\par}
\typeout{TCELL:p173_1:10:\the\ht0:\the\dp0}
\setbox0=\vbox{\hsize=5.440cm\parindent=0pt\linespread{1}\fontsize{10}{13.50}\selectfont\raggedright\noindent 上文“偷救命药案”，甲同时构成盗窃罪和故意杀人罪，想象竞合，择一重罪论处，定故意杀人罪\par}
\typeout{TCELL:p173_1:11:\the\ht0:\the\dp0}
\setbox0=\vbox{\hsize=2.440cm\parindent=0pt\linespread{1}\fontsize{10}{13.50}\selectfont\bfseries\centering\noindent 危害公共安全\par}
\typeout{TCELL:p316_1:0:\the\ht0:\the\dp0}
\setbox0=\vbox{\hsize=3.440cm\parindent=0pt\linespread{1}\fontsize{10}{13.50}\selectfont\bfseries\centering\noindent 要点\par}
\typeout{TCELL:p316_1:1:\the\ht0:\the\dp0}
\setbox0=\vbox{\hsize=9.340cm\parindent=0pt\linespread{1}\fontsize{10}{13.50}\selectfont\bfseries\centering\noindent 范例\par}
\typeout{TCELL:p316_1:2:\the\ht0:\the\dp0}
\setbox0=\vbox{\hsize=2.440cm\parindent=0pt\linespread{1}\fontsize{10}{13.50}\selectfont\centering\noindent 安全\par}
\typeout{TCELL:p316_1:3:\the\ht0:\the\dp0}
\setbox0=\vbox{\hsize=3.440cm\parindent=0pt\linespread{1}\fontsize{10}{13.50}\selectfont\raggedright\noindent 安全内容是指人身伤亡、财产损失（物质性损害），不包括精神损害、心理恐慌（非物质性损害）\par}
\typeout{TCELL:p316_1:4:\the\ht0:\the\dp0}
\setbox0=\vbox{\hsize=9.340cm\parindent=0pt\linespread{1}\fontsize{10}{13.50}\selectfont\raggedright\noindent 甲将面粉邮寄给某单位，谎称是炭疽热病菌粉，引起众人恐慌，乱作一团。甲构成投放虚假危险物质罪。该罪不是危害公共安全犯罪，而是妨害社会管理秩序犯罪，因为引起心理恐慌，不属于危害公共安全\par}
\typeout{TCELL:p316_1:5:\the\ht0:\the\dp0}
\setbox0=\vbox{\hsize=2.440cm\parindent=0pt\linespread{1}\fontsize{10}{13.50}\selectfont\centering\noindent 公共\par}
\typeout{TCELL:p316_1:6:\the\ht0:\the\dp0}
\setbox0=\vbox{\hsize=3.440cm\parindent=0pt\linespread{1}\fontsize{10}{13.50}\selectfont\raggedright\noindent 公共是指公众（多数人或不特定人）。危害公众安全，是指\CJKunderwave{危险具有危及多数人的可能性}\par}
\typeout{TCELL:p316_1:7:\the\ht0:\the\dp0}
\setbox0=\vbox{\hsize=9.340cm\parindent=0pt\linespread{1}\fontsize{10}{13.50}\selectfont\raggedright\noindent 第一，虽然死亡人数有三人，但是危险的危及范围非常确定，只会危及特定少数人，不属于危害公共安全。例如（灭门案），甲向乙家锅里投毒，杀死了乙一家三口，不构成投放危险物质罪，只构成故意杀人罪\par}
\typeout{TCELL:p316_1:8:\the\ht0:\the\dp0}
\setbox0=\vbox{\hsize=9.340cm\parindent=0pt\linespread{1}\fontsize{10}{13.50}\selectfont\raggedright\noindent 第二，虽然死亡人数只有一人，但是危险具有危及多数人的可能性，属于危害公共安全。例如（食堂锅里投毒案），甲是校园食堂师傅，向一锅汤里放老鼠药，一名学生中毒死亡。由于危险的危及范围可大可小，所以甲制造了公共危险，危害了公共安全，构成投放危险物质罪①\par}
\typeout{TCELL:p316_1:9:\the\ht0:\the\dp0}
\setbox0=\vbox{\hsize=2.340cm\parindent=0pt\linespread{1}\fontsize{10}{13.50}\selectfont\bfseries\centering\noindent 妨害公务罪\par}
\typeout{TCELL:p360_1:0:\the\ht0:\the\dp0}
\setbox0=\vbox{\hsize=2.840cm\parindent=0pt\linespread{1}\fontsize{10}{13.50}\selectfont\bfseries\centering\noindent 内容\par}
\typeout{TCELL:p360_1:1:\the\ht0:\the\dp0}
\setbox0=\vbox{\hsize=10.040cm\parindent=0pt\linespread{1}\fontsize{10}{13.50}\selectfont\bfseries\centering\noindent 范例\par}
\typeout{TCELL:p360_1:2:\the\ht0:\the\dp0}
\setbox0=\vbox{\hsize=2.340cm\parindent=0pt\linespread{1}\fontsize{10}{13.50}\selectfont\centering\noindent 妨害的对象\par}
\typeout{TCELL:p360_1:3:\the\ht0:\the\dp0}
\setbox0=\vbox{\hsize=2.840cm\parindent=0pt\linespread{1}\fontsize{10}{13.50}\selectfont\centering\noindent 正在执行的公务\par}
\typeout{TCELL:p360_1:4:\the\ht0:\the\dp0}
\setbox0=\vbox{\hsize=10.040cm\parindent=0pt\linespread{1}\fontsize{10}{13.50}\selectfont\raggedright\noindent 警察调解完甲乙的纠纷，甲乙签字后，甲心里又不平，用脚踹警察，致其轻伤。甲不构成妨害公务罪，而构成故意伤害罪\par}
\typeout{TCELL:p360_1:5:\the\ht0:\the\dp0}
\setbox0=\vbox{\hsize=2.340cm\parindent=0pt\linespread{1}\fontsize{10}{13.50}\selectfont\centering\noindent 认识错误\par}
\typeout{TCELL:p360_1:6:\the\ht0:\the\dp0}
\setbox0=\vbox{\hsize=2.840cm\parindent=0pt\linespread{1}\fontsize{10}{13.50}\selectfont\centering\noindent 故意伤害罪\par}
\typeout{TCELL:p360_1:7:\the\ht0:\the\dp0}
\setbox0=\vbox{\hsize=10.040cm\parindent=0pt\linespread{1}\fontsize{10}{13.50}\selectfont\raggedright\noindent 甲看到朋友乙与丙打架，便上前帮忙，将丙打成轻伤，实际上丙是便衣警察，正在抓捕乙。（1）妨害公务罪：客观要件（$\checkmark$），主观要件（$\times$）。（2）故意伤害罪：客观要件（$\checkmark$），主观要件（$\checkmark$），定故意伤害罪\par}
\typeout{TCELL:p360_1:8:\the\ht0:\the\dp0}
\setbox0=\vbox{\hsize=2.840cm\parindent=0pt\linespread{1}\fontsize{10}{13.50}\selectfont\centering\noindent 假想防卫\par}
\typeout{TCELL:p360_1:9:\the\ht0:\the\dp0}
\setbox0=\vbox{\hsize=10.040cm\parindent=0pt\linespread{1}\fontsize{10}{13.50}\selectfont\raggedright\noindent 甲误以为丙在追杀乙，便将丙打成重伤，实际上丙是便衣警察，正在追捕乙。（1）妨害公务罪：客观要件（$\checkmark$），主观要件（$\times$）。（2）故意伤害罪：客观要件（$\checkmark$），主观要件（$\times$）。如果有过失，则是过失犯罪；如果没有过失，则无罪（2013 年主观题）\par}
\typeout{TCELL:p360_1:10:\the\ht0:\the\dp0}
\setbox0=\vbox{\hsize=2.340cm\parindent=0pt\linespread{1}\fontsize{10}{13.50}\selectfont\centering\noindent 罪数总结\par}
\typeout{TCELL:p360_1:11:\the\ht0:\the\dp0}
\setbox0=\vbox{\hsize=13.140cm\parindent=0pt\linespread{1}\fontsize{10}{13.50}\selectfont\raggedright\noindent 原则上，实施某种犯罪，又实施妨害公务罪，应数罪并罚\par}
\typeout{TCELL:p360_1:12:\the\ht0:\the\dp0}
\setbox0=\vbox{\hsize=13.140cm\parindent=0pt\linespread{1}\fontsize{10}{13.50}\selectfont\raggedright\noindent 例外的，结合犯：前罪+妨害公务罪=前罪，加重处罚。这种前罪只有：（1）走私、贩卖、运输、制造毒品犯罪。（2）组织、运送他人偷越国（边）境罪\par}
\typeout{TCELL:p360_1:13:\the\ht0:\the\dp0}
\setbox0=\vbox{\hsize=15.740cm\parindent=0pt\linespread{1}\fontsize{10}{13.50}\selectfont\bfseries\centering\noindent 贪污罪=A+B\par}
\typeout{TCELL:p397_1:0:\the\ht0:\the\dp0}
\setbox0=\vbox{\hsize=5.040cm\parindent=0pt\linespread{1}\fontsize{10}{13.50}\selectfont\centering\noindent A：实行行为\\（侵占罪、盗窃罪、诈骗罪）\par}
\typeout{TCELL:p397_1:1:\the\ht0:\the\dp0}
\setbox0=\vbox{\hsize=10.440cm\parindent=0pt\linespread{1}\fontsize{10}{13.50}\selectfont\raggedright\noindent B1：行为主体（国家工作人员）\par}
\typeout{TCELL:p397_1:2:\the\ht0:\the\dp0}
\setbox0=\vbox{\hsize=10.440cm\parindent=0pt\linespread{1}\fontsize{10}{13.50}\selectfont\raggedright\noindent B2：行为对象（公共财物，入罪数额是 3 万元）\par}
\typeout{TCELL:p397_1:3:\the\ht0:\the\dp0}
\setbox0=\vbox{\hsize=10.440cm\parindent=0pt\linespread{1}\fontsize{10}{13.50}\selectfont\raggedright\noindent B3：行为方式（利用职务便利）\par}
\typeout{TCELL:p397_1:4:\the\ht0:\the\dp0}
\setbox0=\vbox{\hsize=7.940cm\parindent=0pt\linespread{1}\fontsize{10}{13.50}\selectfont\bfseries\centering\noindent 单位犯罪后单位“没了”\par}
\typeout{TCELL:p041_1:0:\the\ht0:\the\dp0}
\setbox0=\vbox{\hsize=7.540cm\parindent=0pt\linespread{1}\fontsize{10}{13.50}\selectfont\bfseries\centering\noindent 处理结论\par}
\typeout{TCELL:p041_1:1:\the\ht0:\the\dp0}
\setbox0=\vbox{\hsize=1.240cm\parindent=0pt\linespread{1}\fontsize{10}{13.50}\selectfont\centering\noindent 真没了\par}
\typeout{TCELL:p041_1:2:\the\ht0:\the\dp0}
\setbox0=\vbox{\hsize=6.440cm\parindent=0pt\linespread{1}\fontsize{10}{13.50}\selectfont\raggedright\noindent 被撤销、注销、吊销营业执照、宣告破产\par}
\typeout{TCELL:p041_1:3:\the\ht0:\the\dp0}
\setbox0=\vbox{\hsize=7.540cm\parindent=0pt\linespread{1}\fontsize{10}{13.50}\selectfont\raggedright\noindent 直接追究直接责任人的刑事责任（2015 年第 54 题）\par}
\typeout{TCELL:p041_1:4:\the\ht0:\the\dp0}
\setbox0=\vbox{\hsize=1.240cm\parindent=0pt\linespread{1}\fontsize{10}{13.50}\selectfont\centering\noindent 假没了\par}
\typeout{TCELL:p041_1:5:\the\ht0:\the\dp0}
\setbox0=\vbox{\hsize=6.440cm\parindent=0pt\linespread{1}\fontsize{10}{13.50}\selectfont\centering\noindent 被合并到一个新单位\par}
\typeout{TCELL:p041_1:6:\the\ht0:\the\dp0}
\setbox0=\vbox{\hsize=7.540cm\parindent=0pt\linespread{1}\fontsize{10}{13.50}\selectfont\raggedright\noindent 追究原单位的单位犯罪（2018 年试题）\par}
\typeout{TCELL:p041_1:7:\the\ht0:\the\dp0}
\setbox0=\vbox{\hsize=4.340cm\parindent=0pt\linespread{1}\fontsize{10}{13.50}\selectfont\bfseries\centering\noindent 罪过形式\par}
\typeout{TCELL:p072_1:0:\the\ht0:\the\dp0}
\setbox0=\vbox{\hsize=5.540cm\parindent=0pt\linespread{1}\fontsize{10}{13.50}\selectfont\bfseries\centering\noindent 认识因素\par}
\typeout{TCELL:p072_1:1:\the\ht0:\the\dp0}
\setbox0=\vbox{\hsize=5.340cm\parindent=0pt\linespread{1}\fontsize{10}{13.50}\selectfont\bfseries\centering\noindent 意志因素\par}
\typeout{TCELL:p072_1:2:\the\ht0:\the\dp0}
\setbox0=\vbox{\hsize=4.340cm\parindent=0pt\linespread{1}\fontsize{10}{13.50}\selectfont\raggedright\noindent 直接故意\par}
\typeout{TCELL:p072_1:3:\the\ht0:\the\dp0}
\setbox0=\vbox{\hsize=5.540cm\parindent=0pt\linespread{1}\fontsize{10}{13.50}\selectfont\raggedright\noindent 认识到必然或可能发生\par}
\typeout{TCELL:p072_1:4:\the\ht0:\the\dp0}
\setbox0=\vbox{\hsize=5.340cm\parindent=0pt\linespread{1}\fontsize{10}{13.50}\selectfont\raggedright\noindent 希望（赞成票）\par}
\typeout{TCELL:p072_1:5:\the\ht0:\the\dp0}
\setbox0=\vbox{\hsize=4.340cm\parindent=0pt\linespread{1}\fontsize{10}{13.50}\selectfont\raggedright\noindent 间接故意\par}
\typeout{TCELL:p072_1:6:\the\ht0:\the\dp0}
\setbox0=\vbox{\hsize=5.540cm\parindent=0pt\linespread{1}\fontsize{10}{13.50}\selectfont\raggedright\noindent 认识到可能发生\par}
\typeout{TCELL:p072_1:7:\the\ht0:\the\dp0}
\setbox0=\vbox{\hsize=5.340cm\parindent=0pt\linespread{1}\fontsize{10}{13.50}\selectfont\raggedright\noindent 放任（弃权票）\par}
\typeout{TCELL:p072_1:8:\the\ht0:\the\dp0}
\setbox0=\vbox{\hsize=4.340cm\parindent=0pt\linespread{1}\fontsize{10}{13.50}\selectfont\raggedright\noindent 过于自信的过失\par}
\typeout{TCELL:p072_1:9:\the\ht0:\the\dp0}
\setbox0=\vbox{\hsize=5.540cm\parindent=0pt\linespread{1}\fontsize{10}{13.50}\selectfont\raggedright\noindent 认识到可能发生，本应避免\par}
\typeout{TCELL:p072_1:10:\the\ht0:\the\dp0}
\setbox0=\vbox{\hsize=5.340cm\parindent=0pt\linespread{1}\fontsize{10}{13.50}\selectfont\raggedright\noindent 不想发生（反对票）\par}
\typeout{TCELL:p072_1:11:\the\ht0:\the\dp0}
\setbox0=\vbox{\hsize=4.340cm\parindent=0pt\linespread{1}\fontsize{10}{13.50}\selectfont\raggedright\noindent 疏忽大意的过失\par}
\typeout{TCELL:p072_1:12:\the\ht0:\the\dp0}
\setbox0=\vbox{\hsize=5.540cm\parindent=0pt\linespread{1}\fontsize{10}{13.50}\selectfont\raggedright\noindent 没有预见到，但应当预见到\par}
\typeout{TCELL:p072_1:13:\the\ht0:\the\dp0}
\setbox0=\vbox{\hsize=5.340cm\parindent=0pt\linespread{1}\fontsize{10}{13.50}\selectfont\raggedright\noindent 不想发生\par}
\typeout{TCELL:p072_1:14:\the\ht0:\the\dp0}
\setbox0=\vbox{\hsize=4.340cm\parindent=0pt\linespread{1}\fontsize{10}{13.50}\selectfont\raggedright\noindent 意外事件\par}
\typeout{TCELL:p072_1:15:\the\ht0:\the\dp0}
\setbox0=\vbox{\hsize=5.540cm\parindent=0pt\linespread{1}\fontsize{10}{13.50}\selectfont\raggedright\noindent 没有预见到，且无法预见到\par}
\typeout{TCELL:p072_1:16:\the\ht0:\the\dp0}
\setbox0=\vbox{\hsize=5.340cm\parindent=0pt\linespread{1}\fontsize{10}{13.50}\selectfont\raggedright\noindent 不想发生\par}
\typeout{TCELL:p072_1:17:\the\ht0:\the\dp0}
\setbox0=\vbox{\hsize=4.340cm\parindent=0pt\linespread{1}\fontsize{10}{13.50}\selectfont\raggedright\noindent 不可抗力\par}
\typeout{TCELL:p072_1:18:\the\ht0:\the\dp0}
\setbox0=\vbox{\hsize=5.540cm\parindent=0pt\linespread{1}\fontsize{10}{13.50}\selectfont\raggedright\noindent 预见到，但无法避免\par}
\typeout{TCELL:p072_1:19:\the\ht0:\the\dp0}
\setbox0=\vbox{\hsize=5.340cm\parindent=0pt\linespread{1}\fontsize{10}{13.50}\selectfont\raggedright\noindent 不想发生\par}
\typeout{TCELL:p072_1:20:\the\ht0:\the\dp0}
\setbox0=\vbox{\hsize=5.240cm\parindent=0pt\linespread{1}\fontsize{10}{13.50}\selectfont\bfseries\centering\noindent 正当防卫的成立条件\par}
\typeout{TCELL:p088_1:0:\the\ht0:\the\dp0}
\setbox0=\vbox{\hsize=10.240cm\parindent=0pt\linespread{1}\fontsize{10}{13.50}\selectfont\bfseries\centering\noindent 具体内容\par}
\typeout{TCELL:p088_1:1:\the\ht0:\the\dp0}
\setbox0=\vbox{\hsize=5.240cm\parindent=0pt\linespread{1}\fontsize{10}{13.50}\selectfont\raggedright\noindent 起因条件\par}
\typeout{TCELL:p088_1:2:\the\ht0:\the\dp0}
\setbox0=\vbox{\hsize=10.240cm\parindent=0pt\linespread{1}\fontsize{10}{13.50}\selectfont\raggedright\noindent 面临的侵害具有不法性、客观性和现实性（假想防卫）\par}
\typeout{TCELL:p088_1:3:\the\ht0:\the\dp0}
\setbox0=\vbox{\hsize=5.240cm\parindent=0pt\linespread{1}\fontsize{10}{13.50}\selectfont\raggedright\noindent 时间条件\par}
\typeout{TCELL:p088_1:4:\the\ht0:\the\dp0}
\setbox0=\vbox{\hsize=10.240cm\parindent=0pt\linespread{1}\fontsize{10}{13.50}\selectfont\raggedright\noindent 不法侵害正在进行（事后防卫）\par}
\typeout{TCELL:p088_1:5:\the\ht0:\the\dp0}
\setbox0=\vbox{\hsize=5.240cm\parindent=0pt\linespread{1}\fontsize{10}{13.50}\selectfont\raggedright\noindent 意思条件\par}
\typeout{TCELL:p088_1:6:\the\ht0:\the\dp0}
\setbox0=\vbox{\hsize=10.240cm\parindent=0pt\linespread{1}\fontsize{10}{13.50}\selectfont\raggedright\noindent 防卫者是否需要防卫意思（偶然防卫）\par}
\typeout{TCELL:p088_1:7:\the\ht0:\the\dp0}
\setbox0=\vbox{\hsize=5.240cm\parindent=0pt\linespread{1}\fontsize{10}{13.50}\selectfont\raggedright\noindent 对象条件\par}
\typeout{TCELL:p088_1:8:\the\ht0:\the\dp0}
\setbox0=\vbox{\hsize=10.240cm\parindent=0pt\linespread{1}\fontsize{10}{13.50}\selectfont\raggedright\noindent 防卫手段针对不法侵害人本人\par}
\typeout{TCELL:p088_1:9:\the\ht0:\the\dp0}
\setbox0=\vbox{\hsize=5.240cm\parindent=0pt\linespread{1}\fontsize{10}{13.50}\selectfont\raggedright\noindent 限度条件\par}
\typeout{TCELL:p088_1:10:\the\ht0:\the\dp0}
\setbox0=\vbox{\hsize=10.240cm\parindent=0pt\linespread{1}\fontsize{10}{13.50}\selectfont\raggedright\noindent 防卫手段具有必要性和相当性（防卫过当）\par}
\typeout{TCELL:p088_1:11:\the\ht0:\the\dp0}
\setbox0=\vbox{\hsize=3.140cm\parindent=0pt\linespread{1}\fontsize{10}{13.50}\selectfont\centering\noindent 客观要件\par}
\typeout{TCELL:p090_1:0:\the\ht0:\the\dp0}
\setbox0=\vbox{\hsize=5.540cm\parindent=0pt\linespread{1}\fontsize{10}{13.50}\selectfont\centering\noindent （作为或不作为的）不法侵害\par}
\typeout{TCELL:p090_1:1:\the\ht0:\the\dp0}
\setbox0=\vbox{\hsize=3.140cm\parindent=0pt\linespread{1}\fontsize{10}{13.50}\selectfont\centering\noindent 主观要件\par}
\typeout{TCELL:p090_1:2:\the\ht0:\the\dp0}
\setbox0=\vbox{\hsize=5.540cm\parindent=0pt\linespread{1}\fontsize{10}{13.50}\selectfont\centering\noindent 故意的不法侵害\par}
\typeout{TCELL:p090_1:3:\the\ht0:\the\dp0}
\setbox0=\vbox{\hsize=5.540cm\parindent=0pt\linespread{1}\fontsize{10}{13.50}\selectfont\centering\noindent 过失的不法侵害\par}
\typeout{TCELL:p090_1:4:\the\ht0:\the\dp0}
\setbox0=\vbox{\hsize=3.140cm\parindent=0pt\linespread{1}\fontsize{10}{13.50}\selectfont\centering\noindent 排除责任事由\par}
\typeout{TCELL:p090_1:5:\the\ht0:\the\dp0}
\setbox0=\vbox{\hsize=5.540cm\parindent=0pt\linespread{1}\fontsize{10}{13.50}\selectfont\centering\noindent 无责任年龄的不法侵害\par}
\typeout{TCELL:p090_1:6:\the\ht0:\the\dp0}
\setbox0=\vbox{\hsize=5.540cm\parindent=0pt\linespread{1}\fontsize{10}{13.50}\selectfont\centering\noindent 无责任能力的不法侵害\par}
\typeout{TCELL:p090_1:7:\the\ht0:\the\dp0}
\setbox0=\vbox{\hsize=5.240cm\parindent=0pt\linespread{1}\fontsize{10}{13.50}\selectfont\bfseries\centering\noindent 制止不法侵害\par}
\typeout{TCELL:p094_2:0:\the\ht0:\the\dp0}
\setbox0=\vbox{\hsize=5.240cm\parindent=0pt\linespread{1}\fontsize{10}{13.50}\selectfont\bfseries\centering\noindent 偶然防卫（坏心办好事）\par}
\typeout{TCELL:p094_2:1:\the\ht0:\the\dp0}
\setbox0=\vbox{\hsize=4.740cm\parindent=0pt\linespread{1}\fontsize{10}{13.50}\selectfont\bfseries\centering\noindent 假想防卫（好心办坏事）\par}
\typeout{TCELL:p094_2:2:\the\ht0:\the\dp0}
\setbox0=\vbox{\hsize=5.240cm\parindent=0pt\linespread{1}\fontsize{10}{13.50}\selectfont\raggedright\noindent 客观要件\par}
\typeout{TCELL:p094_2:3:\the\ht0:\the\dp0}
\setbox0=\vbox{\hsize=5.240cm\parindent=0pt\linespread{1}\fontsize{10}{13.50}\selectfont\raggedright\noindent $\checkmark$\par}
\typeout{TCELL:p094_2:4:\the\ht0:\the\dp0}
\setbox0=\vbox{\hsize=4.740cm\parindent=0pt\linespread{1}\fontsize{10}{13.50}\selectfont\raggedright\noindent $\times$\par}
\typeout{TCELL:p094_2:5:\the\ht0:\the\dp0}
\setbox0=\vbox{\hsize=5.240cm\parindent=0pt\linespread{1}\fontsize{10}{13.50}\selectfont\raggedright\noindent 主观要件\par}
\typeout{TCELL:p094_2:6:\the\ht0:\the\dp0}
\setbox0=\vbox{\hsize=5.240cm\parindent=0pt\linespread{1}\fontsize{10}{13.50}\selectfont\raggedright\noindent $\times$\par}
\typeout{TCELL:p094_2:7:\the\ht0:\the\dp0}
\setbox0=\vbox{\hsize=4.740cm\parindent=0pt\linespread{1}\fontsize{10}{13.50}\selectfont\raggedright\noindent $\checkmark$\par}
\typeout{TCELL:p094_2:8:\the\ht0:\the\dp0}
\setbox0=\vbox{\hsize=5.240cm\parindent=0pt\linespread{1}\fontsize{10}{13.50}\selectfont\raggedright\noindent 结论\par}
\typeout{TCELL:p094_2:9:\the\ht0:\the\dp0}
\setbox0=\vbox{\hsize=5.240cm\parindent=0pt\linespread{1}\fontsize{10}{13.50}\selectfont\raggedright\noindent 观点展示\par}
\typeout{TCELL:p094_2:10:\the\ht0:\the\dp0}
\setbox0=\vbox{\hsize=4.740cm\parindent=0pt\linespread{1}\fontsize{10}{13.50}\selectfont\raggedright\noindent 不构成正当防卫\par}
\typeout{TCELL:p094_2:11:\the\ht0:\the\dp0}
\setbox0=\vbox{\hsize=3.740cm\parindent=0pt\linespread{1}\fontsize{10}{13.50}\selectfont\bfseries\centering\noindent 偶然避险的辨认标准\par}
\typeout{TCELL:p101_1:0:\the\ht0:\the\dp0}
\setbox0=\vbox{\hsize=2.940cm\parindent=0pt\linespread{1}\fontsize{10}{13.50}\selectfont\bfseries\centering\noindent 范例\par}
\typeout{TCELL:p101_1:1:\the\ht0:\the\dp0}
\setbox0=\vbox{\hsize=4.340cm\parindent=0pt\linespread{1}\fontsize{10}{13.50}\selectfont\bfseries\centering\noindent 观点展示\par}
\typeout{TCELL:p101_1:2:\the\ht0:\the\dp0}
\setbox0=\vbox{\hsize=3.940cm\parindent=0pt\linespread{1}\fontsize{10}{13.50}\selectfont\bfseries\centering\noindent 处理结论\par}
\typeout{TCELL:p101_1:3:\the\ht0:\the\dp0}
\setbox0=\vbox{\hsize=3.740cm\parindent=0pt\linespread{1}\fontsize{10}{13.50}\selectfont\raggedright\noindent 在避免危险这一点上，没有做到主客观相一致，也即客观上避免了危险，但行为人主观上没有认识到这一点\par}
\typeout{TCELL:p101_1:4:\the\ht0:\the\dp0}
\setbox0=\vbox{\hsize=2.940cm\parindent=0pt\linespread{1}\fontsize{10}{13.50}\selectfont\raggedright\noindent 甲以毁坏财物的目的砸坏乙的汽车玻璃，却碰巧救了被锁在车里的快要窒息的小孩\par}
\typeout{TCELL:p101_1:5:\the\ht0:\the\dp0}
\setbox0=\vbox{\hsize=4.340cm\parindent=0pt\linespread{1}\fontsize{10}{13.50}\selectfont\raggedright\noindent 避险认识不要说：成立紧急避险，只要求客观上避免了危险，不要求主观上认识到这一点。结论：偶然避险成立紧急避险\par}
\typeout{TCELL:p101_1:6:\the\ht0:\the\dp0}
\setbox0=\vbox{\hsize=3.940cm\parindent=0pt\linespread{1}\fontsize{10}{13.50}\selectfont\raggedright\noindent 甲成立紧急避险，不构成故意毁坏财物罪\par}
\typeout{TCELL:p101_1:7:\the\ht0:\the\dp0}
\setbox0=\vbox{\hsize=4.340cm\parindent=0pt\linespread{1}\fontsize{10}{13.50}\selectfont\raggedright\noindent 避险认识必要说：成立紧急避险，要求主观上具有避险认识。结论：偶然避险不成立紧急避险\par}
\typeout{TCELL:p101_1:8:\the\ht0:\the\dp0}
\setbox0=\vbox{\hsize=3.940cm\parindent=0pt\linespread{1}\fontsize{10}{13.50}\selectfont\raggedright\noindent 甲不成立紧急避险，成立故意毁坏财物罪，但是由于甲制造了好结果，所以不构成故意毁坏财物罪既遂，而构成未遂\par}
\typeout{TCELL:p101_1:9:\the\ht0:\the\dp0}
\setbox0=\vbox{\hsize=1.040cm\parindent=0pt\linespread{1}\fontsize{10}{13.50}\selectfont\bfseries\centering\noindent 种类\par}
\typeout{TCELL:p102_1:0:\the\ht0:\the\dp0}
\setbox0=\vbox{\hsize=3.540cm\parindent=0pt\linespread{1}\fontsize{10}{13.50}\selectfont\bfseries\centering\noindent 起因条件\par}
\typeout{TCELL:p102_1:1:\the\ht0:\the\dp0}
\setbox0=\vbox{\hsize=2.540cm\parindent=0pt\linespread{1}\fontsize{10}{13.50}\selectfont\bfseries\centering\noindent 时间条件\par}
\typeout{TCELL:p102_1:2:\the\ht0:\the\dp0}
\setbox0=\vbox{\hsize=2.440cm\parindent=0pt\linespread{1}\fontsize{10}{13.50}\selectfont\bfseries\centering\noindent 意思条件\par}
\typeout{TCELL:p102_1:3:\the\ht0:\the\dp0}
\setbox0=\vbox{\hsize=2.440cm\parindent=0pt\linespread{1}\fontsize{10}{13.50}\selectfont\bfseries\centering\noindent “不得已”条件\par}
\typeout{TCELL:p102_1:4:\the\ht0:\the\dp0}
\setbox0=\vbox{\hsize=2.440cm\parindent=0pt\linespread{1}\fontsize{10}{13.50}\selectfont\bfseries\centering\noindent 限度条件\par}
\typeout{TCELL:p102_1:5:\the\ht0:\the\dp0}
\setbox0=\vbox{\hsize=1.040cm\parindent=0pt\linespread{1}\fontsize{10}{13.50}\selectfont\raggedright\noindent 正当防卫\par}
\typeout{TCELL:p102_1:6:\the\ht0:\the\dp0}
\setbox0=\vbox{\hsize=3.540cm\parindent=0pt\linespread{1}\fontsize{10}{13.50}\selectfont\raggedright\noindent 制止不法侵害。可以为了保护国家法益、公共法益而实施正当防卫\par}
\typeout{TCELL:p102_1:7:\the\ht0:\the\dp0}
\setbox0=\vbox{\hsize=2.540cm\parindent=0pt\linespread{1}\fontsize{10}{13.50}\selectfont\raggedright\noindent 不法侵害正在进行，紧迫程度高一些\par}
\typeout{TCELL:p102_1:8:\the\ht0:\the\dp0}
\setbox0=\vbox{\hsize=2.440cm\parindent=0pt\linespread{1}\fontsize{10}{13.50}\selectfont\raggedright\noindent 偶然防卫，防卫认识不要说与必要说\par}
\typeout{TCELL:p102_1:9:\the\ht0:\the\dp0}
\setbox0=\vbox{\hsize=2.440cm\parindent=0pt\linespread{1}\fontsize{10}{13.50}\selectfont\raggedright\noindent 原则上，没有躲避义务，防卫反击是优先手段，不要求不得已而为之\par}
\typeout{TCELL:p102_1:10:\the\ht0:\the\dp0}
\setbox0=\vbox{\hsize=2.440cm\parindent=0pt\linespread{1}\fontsize{10}{13.50}\selectfont\raggedright\noindent 更宽松。防卫明显超过必要限度造成重大损害的，才属于防卫过当\par}
\typeout{TCELL:p102_1:11:\the\ht0:\the\dp0}
\setbox0=\vbox{\hsize=1.040cm\parindent=0pt\linespread{1}\fontsize{10}{13.50}\selectfont\bfseries\centering\noindent 种类\par}
\typeout{TCELL:p103_1:0:\the\ht0:\the\dp0}
\setbox0=\vbox{\hsize=3.540cm\parindent=0pt\linespread{1}\fontsize{10}{13.50}\selectfont\bfseries\centering\noindent 起因条件\par}
\typeout{TCELL:p103_1:1:\the\ht0:\the\dp0}
\setbox0=\vbox{\hsize=2.540cm\parindent=0pt\linespread{1}\fontsize{10}{13.50}\selectfont\bfseries\centering\noindent 时间条件\par}
\typeout{TCELL:p103_1:2:\the\ht0:\the\dp0}
\setbox0=\vbox{\hsize=2.440cm\parindent=0pt\linespread{1}\fontsize{10}{13.50}\selectfont\bfseries\centering\noindent 意思条件\par}
\typeout{TCELL:p103_1:3:\the\ht0:\the\dp0}
\setbox0=\vbox{\hsize=2.440cm\parindent=0pt\linespread{1}\fontsize{10}{13.50}\selectfont\bfseries\centering\noindent “不得已”条件\par}
\typeout{TCELL:p103_1:4:\the\ht0:\the\dp0}
\setbox0=\vbox{\hsize=2.440cm\parindent=0pt\linespread{1}\fontsize{10}{13.50}\selectfont\bfseries\centering\noindent 限度条件\par}
\typeout{TCELL:p103_1:5:\the\ht0:\the\dp0}
\setbox0=\vbox{\hsize=1.040cm\parindent=0pt\linespread{1}\fontsize{10}{13.50}\selectfont\raggedright\noindent 紧急避险\par}
\typeout{TCELL:p103_1:6:\the\ht0:\the\dp0}
\setbox0=\vbox{\hsize=3.540cm\parindent=0pt\linespread{1}\fontsize{10}{13.50}\selectfont\raggedright\noindent 避免危险。可以为了保护国家法益、公共法益而实施紧急避险（2024 年试题）\par}
\typeout{TCELL:p103_1:7:\the\ht0:\the\dp0}
\setbox0=\vbox{\hsize=2.540cm\parindent=0pt\linespread{1}\fontsize{10}{13.50}\selectfont\raggedright\noindent 危险正在发生，紧迫程度低一些\par}
\typeout{TCELL:p103_1:8:\the\ht0:\the\dp0}
\setbox0=\vbox{\hsize=2.440cm\parindent=0pt\linespread{1}\fontsize{10}{13.50}\selectfont\raggedright\noindent 偶然避险，避险认识不要说与必要说\par}
\typeout{TCELL:p103_1:9:\the\ht0:\the\dp0}
\setbox0=\vbox{\hsize=2.440cm\parindent=0pt\linespread{1}\fontsize{10}{13.50}\selectfont\raggedright\noindent 是补充性手段，要求不得已而为之（2024 年试题）\par}
\typeout{TCELL:p103_1:10:\the\ht0:\the\dp0}
\setbox0=\vbox{\hsize=2.440cm\parindent=0pt\linespread{1}\fontsize{10}{13.50}\selectfont\raggedright\noindent 更严格。避险超过必要限度造成不应有的损害的，便属于避险过当（2017 年第 4 题）\par}
\typeout{TCELL:p103_1:11:\the\ht0:\the\dp0}
\setbox0=\vbox{\hsize=2.840cm\parindent=0pt\linespread{1}\fontsize{10}{13.50}\selectfont\bfseries\centering\noindent 反击对象\par}
\typeout{TCELL:p103_2:0:\the\ht0:\the\dp0}
\setbox0=\vbox{\hsize=6.340cm\parindent=0pt\linespread{1}\fontsize{10}{13.50}\selectfont\bfseries\centering\noindent 结论和理由\par}
\typeout{TCELL:p103_2:1:\the\ht0:\the\dp0}
\setbox0=\vbox{\hsize=6.040cm\parindent=0pt\linespread{1}\fontsize{10}{13.50}\selectfont\bfseries\centering\noindent 范例\par}
\typeout{TCELL:p103_2:2:\the\ht0:\the\dp0}
\setbox0=\vbox{\hsize=2.840cm\parindent=0pt\linespread{1}\fontsize{10}{13.50}\selectfont\raggedright\noindent 正当防卫\par}
\typeout{TCELL:p103_2:3:\the\ht0:\the\dp0}
\setbox0=\vbox{\hsize=6.340cm\parindent=0pt\linespread{1}\fontsize{10}{13.50}\selectfont\raggedright\noindent 对正当防卫的反击，不构成正当防卫，而是故意犯罪，因为正当防卫不属于不法侵害\par}
\typeout{TCELL:p103_2:4:\the\ht0:\the\dp0}
\setbox0=\vbox{\hsize=6.040cm\parindent=0pt\linespread{1}\fontsize{10}{13.50}\selectfont\raggedright\noindent 甲杀乙，乙反击甲，甲又反击乙。甲的反击构成故意犯罪\par}
\typeout{TCELL:p103_2:5:\the\ht0:\the\dp0}
\setbox0=\vbox{\hsize=2.840cm\parindent=0pt\linespread{1}\fontsize{10}{13.50}\selectfont\raggedright\noindent 紧急避险\par}
\typeout{TCELL:p103_2:6:\the\ht0:\the\dp0}
\setbox0=\vbox{\hsize=6.340cm\parindent=0pt\linespread{1}\fontsize{10}{13.50}\selectfont\raggedright\noindent （1）对紧急避险的反击，不构成正当防卫，因为紧急避险不属于不法侵害。（2）对紧急避险的反击不构成紧急避险，因为不符合法益衡量的要求。（3）对紧急避险的反击属于不法侵害\par}
\typeout{TCELL:p103_2:7:\the\ht0:\the\dp0}
\setbox0=\vbox{\hsize=6.040cm\parindent=0pt\linespread{1}\fontsize{10}{13.50}\selectfont\raggedright\noindent 几条恶犬追咬甲，甲不得已闯入乙家。乙见状，将甲推出去。乙不构成紧急避险，因为乙保护的法益（住宅安宁状态）小于损害的法益（甲的身体健康）（2024 年试题）\par}
\typeout{TCELL:p103_2:8:\the\ht0:\the\dp0}
\setbox0=\vbox{\hsize=2.840cm\parindent=0pt\linespread{1}\fontsize{10}{13.50}\selectfont\raggedright\noindent 假想防卫、事后防卫、防卫过当；假想避险、事后避险、避险过当\par}
\typeout{TCELL:p103_2:9:\the\ht0:\the\dp0}
\setbox0=\vbox{\hsize=6.340cm\parindent=0pt\linespread{1}\fontsize{10}{13.50}\selectfont\raggedright\noindent 对这些行为的反击，构成正当防卫，因为这些行为都不构成正当防卫、紧急避险，都属于不法侵害\par}
\typeout{TCELL:p103_2:10:\the\ht0:\the\dp0}
\setbox0=\vbox{\hsize=6.040cm\parindent=0pt\linespread{1}\fontsize{10}{13.50}\selectfont\raggedright\noindent 甲伤害乙后，离去，乙追上甲，用刀砍甲。甲反击乙。乙砍甲，属于事后防卫，属于不法侵害。甲反击乙构成正当防卫\par}
\typeout{TCELL:p103_2:11:\the\ht0:\the\dp0}
\setbox0=\vbox{\hsize=2.840cm\parindent=0pt\linespread{1}\fontsize{10}{13.50}\selectfont\raggedright\noindent 偶然防卫、偶然避险\par}
\typeout{TCELL:p103_2:12:\the\ht0:\the\dp0}
\setbox0=\vbox{\hsize=6.340cm\parindent=0pt\linespread{1}\fontsize{10}{13.50}\selectfont\raggedright\noindent 对其反击，观点展示：（1）按照防卫（避险）认识不要说，偶然防卫（偶然避险）构成正当防卫（紧急避险），不是不法侵害，因此对其不能进行正当防卫。（2）按照防卫（避险）认识必要说，偶然防卫（偶然避险）不构成正当防卫（紧急避险），客观上属于不法侵害，因此对其可以进行正当防卫\par}
\typeout{TCELL:p103_2:13:\the\ht0:\the\dp0}
\setbox0=\vbox{\hsize=6.040cm\parindent=0pt\linespread{1}\fontsize{10}{13.50}\selectfont\raggedright\noindent 乙正要开枪杀丙，甲正要开枪杀乙，甲没有看到乙要杀人，甲只想杀死乙。甲的行为属于偶然防卫。丁反击甲。防卫认识不要说认为，丁不构成正当防卫。防卫认识必要说认为，丁构成正当防卫\par}
\typeout{TCELL:p103_2:14:\the\ht0:\the\dp0}
\setbox0=\vbox{\hsize=4.940cm\parindent=0pt\linespread{1}\fontsize{10}{13.50}\selectfont\bfseries\centering\noindent 责任能力的程度\par}
\typeout{TCELL:p109_1:0:\the\ht0:\the\dp0}
\setbox0=\vbox{\hsize=4.940cm\parindent=0pt\linespread{1}\fontsize{10}{13.50}\selectfont\bfseries\centering\noindent 精神状态\par}
\typeout{TCELL:p109_1:1:\the\ht0:\the\dp0}
\setbox0=\vbox{\hsize=5.340cm\parindent=0pt\linespread{1}\fontsize{10}{13.50}\selectfont\bfseries\centering\noindent 处理结论\par}
\typeout{TCELL:p109_1:2:\the\ht0:\the\dp0}
\setbox0=\vbox{\hsize=4.940cm\parindent=0pt\linespread{1}\fontsize{10}{13.50}\selectfont\raggedright\noindent 完全责任能力\par}
\typeout{TCELL:p109_1:3:\the\ht0:\the\dp0}
\setbox0=\vbox{\hsize=4.940cm\parindent=0pt\linespread{1}\fontsize{10}{13.50}\selectfont\raggedright\noindent 完全有辨认和控制能力\par}
\typeout{TCELL:p109_1:4:\the\ht0:\the\dp0}
\setbox0=\vbox{\hsize=5.340cm\parindent=0pt\linespread{1}\fontsize{10}{13.50}\selectfont\raggedright\noindent 应负刑事责任\par}
\typeout{TCELL:p109_1:5:\the\ht0:\the\dp0}
\setbox0=\vbox{\hsize=4.940cm\parindent=0pt\linespread{1}\fontsize{10}{13.50}\selectfont\raggedright\noindent 相对责任能力（又称限定或限制责任能力）\par}
\typeout{TCELL:p109_1:6:\the\ht0:\the\dp0}
\setbox0=\vbox{\hsize=4.940cm\parindent=0pt\linespread{1}\fontsize{10}{13.50}\selectfont\raggedright\noindent 尚未完全丧失辨认和控制能力\par}
\typeout{TCELL:p109_1:7:\the\ht0:\the\dp0}
\setbox0=\vbox{\hsize=5.340cm\parindent=0pt\linespread{1}\fontsize{10}{13.50}\selectfont\raggedright\noindent 应负刑事责任，可以从宽处罚\par}
\typeout{TCELL:p109_1:8:\the\ht0:\the\dp0}
\setbox0=\vbox{\hsize=4.940cm\parindent=0pt\linespread{1}\fontsize{10}{13.50}\selectfont\bfseries\centering\noindent 责任能力的程度\par}
\typeout{TCELL:p110_1:0:\the\ht0:\the\dp0}
\setbox0=\vbox{\hsize=4.940cm\parindent=0pt\linespread{1}\fontsize{10}{13.50}\selectfont\bfseries\centering\noindent 精神状态\par}
\typeout{TCELL:p110_1:1:\the\ht0:\the\dp0}
\setbox0=\vbox{\hsize=5.340cm\parindent=0pt\linespread{1}\fontsize{10}{13.50}\selectfont\bfseries\centering\noindent 处理结论\par}
\typeout{TCELL:p110_1:2:\the\ht0:\the\dp0}
\setbox0=\vbox{\hsize=4.940cm\parindent=0pt\linespread{1}\fontsize{10}{13.50}\selectfont\raggedright\noindent 完全无责任能力\par}
\typeout{TCELL:p110_1:3:\the\ht0:\the\dp0}
\setbox0=\vbox{\hsize=4.940cm\parindent=0pt\linespread{1}\fontsize{10}{13.50}\selectfont\raggedright\noindent 完全丧失辨认和控制能力\par}
\typeout{TCELL:p110_1:4:\the\ht0:\the\dp0}
\setbox0=\vbox{\hsize=5.340cm\parindent=0pt\linespread{1}\fontsize{10}{13.50}\selectfont\raggedright\noindent 不负刑事责任，应当责令他的家属或者监护人严加看管和医疗；在必要的时候，由政府强制医疗\par}
\typeout{TCELL:p110_1:5:\the\ht0:\the\dp0}
\setbox0=\vbox{\hsize=4.340cm\parindent=0pt\linespread{1}\fontsize{10}{13.50}\selectfont\bfseries\centering\noindent 种类\par}
\typeout{TCELL:p115_1:0:\the\ht0:\the\dp0}
\setbox0=\vbox{\hsize=4.440cm\parindent=0pt\linespread{1}\fontsize{10}{13.50}\selectfont\bfseries\centering\noindent 认识对象\par}
\typeout{TCELL:p115_1:1:\the\ht0:\the\dp0}
\setbox0=\vbox{\hsize=6.440cm\parindent=0pt\linespread{1}\fontsize{10}{13.50}\selectfont\bfseries\centering\noindent 法律后果\par}
\typeout{TCELL:p115_1:2:\the\ht0:\the\dp0}
\setbox0=\vbox{\hsize=4.340cm\parindent=0pt\linespread{1}\fontsize{10}{13.50}\selectfont\raggedright\noindent 事实认识错误\par}
\typeout{TCELL:p115_1:3:\the\ht0:\the\dp0}
\setbox0=\vbox{\hsize=4.440cm\parindent=0pt\linespread{1}\fontsize{10}{13.50}\selectfont\raggedright\noindent 构成要件事实（例如行为对象）\par}
\typeout{TCELL:p115_1:4:\the\ht0:\the\dp0}
\setbox0=\vbox{\hsize=6.440cm\parindent=0pt\linespread{1}\fontsize{10}{13.50}\selectfont\raggedright\noindent 排除犯罪故意\par}
\typeout{TCELL:p115_1:5:\the\ht0:\the\dp0}
\setbox0=\vbox{\hsize=4.340cm\parindent=0pt\linespread{1}\fontsize{10}{13.50}\selectfont\raggedright\noindent 法律认识错误\par}
\typeout{TCELL:p115_1:6:\the\ht0:\the\dp0}
\setbox0=\vbox{\hsize=4.440cm\parindent=0pt\linespread{1}\fontsize{10}{13.50}\selectfont\raggedright\noindent 刑法的禁止性（违法性）\par}
\typeout{TCELL:p115_1:7:\the\ht0:\the\dp0}
\setbox0=\vbox{\hsize=6.440cm\parindent=0pt\linespread{1}\fontsize{10}{13.50}\selectfont\raggedright\noindent 若无违法性认识可能性，则排除责任\par}
\typeout{TCELL:p115_1:8:\the\ht0:\the\dp0}
\setbox0=\vbox{\hsize=4.140cm\parindent=0pt\linespread{1}\fontsize{10}{13.50}\selectfont\bfseries\centering\noindent 要件与种类\par}
\typeout{TCELL:p120_1:0:\the\ht0:\the\dp0}
\setbox0=\vbox{\hsize=5.840cm\parindent=0pt\linespread{1}\fontsize{10}{13.50}\selectfont\bfseries\centering\noindent 不能犯\par}
\typeout{TCELL:p120_1:1:\the\ht0:\the\dp0}
\setbox0=\vbox{\hsize=5.240cm\parindent=0pt\linespread{1}\fontsize{10}{13.50}\selectfont\bfseries\centering\noindent 未遂犯\par}
\typeout{TCELL:p120_1:2:\the\ht0:\the\dp0}
\setbox0=\vbox{\hsize=4.140cm\parindent=0pt\linespread{1}\fontsize{10}{13.50}\selectfont\raggedright\noindent 客观要件\par}
\typeout{TCELL:p120_1:3:\the\ht0:\the\dp0}
\setbox0=\vbox{\hsize=5.840cm\parindent=0pt\linespread{1}\fontsize{10}{13.50}\selectfont\raggedright\noindent 危害行为$\times$，因为对法益没有危险\par}
\typeout{TCELL:p120_1:4:\the\ht0:\the\dp0}
\setbox0=\vbox{\hsize=5.240cm\parindent=0pt\linespread{1}\fontsize{10}{13.50}\selectfont\raggedright\noindent 危害行为$\checkmark$，因为对法益有危险\par}
\typeout{TCELL:p120_1:5:\the\ht0:\the\dp0}
\setbox0=\vbox{\hsize=4.140cm\parindent=0pt\linespread{1}\fontsize{10}{13.50}\selectfont\raggedright\noindent 主观要件\par}
\typeout{TCELL:p120_1:6:\the\ht0:\the\dp0}
\setbox0=\vbox{\hsize=5.840cm\parindent=0pt\linespread{1}\fontsize{10}{13.50}\selectfont\raggedright\noindent 犯罪故意$\checkmark$\par}
\typeout{TCELL:p120_1:7:\the\ht0:\the\dp0}
\setbox0=\vbox{\hsize=5.240cm\parindent=0pt\linespread{1}\fontsize{10}{13.50}\selectfont\raggedright\noindent 犯罪故意$\checkmark$\par}
\typeout{TCELL:p120_1:8:\the\ht0:\the\dp0}
\setbox0=\vbox{\hsize=4.140cm\parindent=0pt\linespread{1}\fontsize{10}{13.50}\selectfont\raggedright\noindent 结论\par}
\typeout{TCELL:p120_1:9:\the\ht0:\the\dp0}
\setbox0=\vbox{\hsize=5.840cm\parindent=0pt\linespread{1}\fontsize{10}{13.50}\selectfont\raggedright\noindent 无罪\par}
\typeout{TCELL:p120_1:10:\the\ht0:\the\dp0}
\setbox0=\vbox{\hsize=5.240cm\parindent=0pt\linespread{1}\fontsize{10}{13.50}\selectfont\raggedright\noindent 有罪，未遂\par}
\typeout{TCELL:p120_1:11:\the\ht0:\the\dp0}
\setbox0=\vbox{\hsize=5.840cm\parindent=0pt\linespread{1}\fontsize{10}{13.50}\selectfont\bfseries\centering\noindent 案件事实特征\par}
\typeout{TCELL:p121_1:0:\the\ht0:\the\dp0}
\setbox0=\vbox{\hsize=0.840cm\parindent=0pt\linespread{1}\fontsize{10}{13.50}\selectfont\bfseries\centering\noindent 结论\par}
\typeout{TCELL:p121_1:1:\the\ht0:\the\dp0}
\setbox0=\vbox{\hsize=8.540cm\parindent=0pt\linespread{1}\fontsize{10}{13.50}\selectfont\bfseries\centering\noindent 名称（术语）\par}
\typeout{TCELL:p121_1:2:\the\ht0:\the\dp0}
\setbox0=\vbox{\hsize=4.640cm\parindent=0pt\linespread{1}\fontsize{10}{13.50}\selectfont\bfseries\centering\noindent 简洁名称（张明楷教授的观点）\par}
\typeout{TCELL:p121_1:3:\the\ht0:\the\dp0}
\setbox0=\vbox{\hsize=3.640cm\parindent=0pt\linespread{1}\fontsize{10}{13.50}\selectfont\bfseries\centering\noindent 复杂名称（传统观点）\par}
\typeout{TCELL:p121_1:4:\the\ht0:\the\dp0}
\setbox0=\vbox{\hsize=5.840cm\parindent=0pt\linespread{1}\fontsize{10}{13.50}\selectfont\raggedright\noindent 行为对法益没有危险（例如，在荒无人烟的沙漠，甲误以为稻草人是仇人而开枪）\par}
\typeout{TCELL:p121_1:5:\the\ht0:\the\dp0}
\setbox0=\vbox{\hsize=0.840cm\parindent=0pt\linespread{1}\fontsize{10}{13.50}\selectfont\centering\noindent 无罪\par}
\typeout{TCELL:p121_1:6:\the\ht0:\the\dp0}
\setbox0=\vbox{\hsize=4.640cm\parindent=0pt\linespread{1}\fontsize{10}{13.50}\selectfont\centering\noindent 不能犯\par}
\typeout{TCELL:p121_1:7:\the\ht0:\the\dp0}
\setbox0=\vbox{\hsize=3.640cm\parindent=0pt\linespread{1}\fontsize{10}{13.50}\selectfont\raggedright\noindent 绝对不能犯，不可罚的不能犯\par}
\typeout{TCELL:p121_1:8:\the\ht0:\the\dp0}
\setbox0=\vbox{\hsize=5.840cm\parindent=0pt\linespread{1}\fontsize{10}{13.50}\selectfont\raggedright\noindent 行为对法益有危险（例如，甲误将砒霜旁边的白糖当作砒霜，投给乙）\par}
\typeout{TCELL:p121_1:9:\the\ht0:\the\dp0}
\setbox0=\vbox{\hsize=0.840cm\parindent=0pt\linespread{1}\fontsize{10}{13.50}\selectfont\centering\noindent 有罪\par}
\typeout{TCELL:p121_1:10:\the\ht0:\the\dp0}
\setbox0=\vbox{\hsize=4.640cm\parindent=0pt\linespread{1}\fontsize{10}{13.50}\selectfont\centering\noindent 未遂犯\par}
\typeout{TCELL:p121_1:11:\the\ht0:\the\dp0}
\setbox0=\vbox{\hsize=3.640cm\parindent=0pt\linespread{1}\fontsize{10}{13.50}\selectfont\raggedright\noindent 相对不能犯，可罚的不能犯，不能犯未遂\par}
\typeout{TCELL:p121_1:12:\the\ht0:\the\dp0}
\setbox0=\vbox{\hsize=2.040cm\parindent=0pt\linespread{1}\fontsize{10}{13.50}\selectfont\bfseries\centering\noindent 可能有效性\par}
\typeout{TCELL:p127_1:0:\the\ht0:\the\dp0}
\setbox0=\vbox{\hsize=3.040cm\parindent=0pt\linespread{1}\fontsize{10}{13.50}\selectfont\bfseries\centering\noindent 实际有效性\par}
\typeout{TCELL:p127_1:1:\the\ht0:\the\dp0}
\setbox0=\vbox{\hsize=1.040cm\parindent=0pt\linespread{1}\fontsize{10}{13.50}\selectfont\bfseries\centering\noindent 结论\par}
\typeout{TCELL:p127_1:2:\the\ht0:\the\dp0}
\setbox0=\vbox{\hsize=8.840cm\parindent=0pt\linespread{1}\fontsize{10}{13.50}\selectfont\bfseries\centering\noindent 范例\par}
\typeout{TCELL:p127_1:3:\the\ht0:\the\dp0}
\setbox0=\vbox{\hsize=2.040cm\parindent=0pt\linespread{1}\fontsize{10}{13.50}\selectfont\raggedright\noindent $\checkmark$\par}
\typeout{TCELL:p127_1:4:\the\ht0:\the\dp0}
\setbox0=\vbox{\hsize=3.040cm\parindent=0pt\linespread{1}\fontsize{10}{13.50}\selectfont\raggedright\noindent $\checkmark$\par}
\typeout{TCELL:p127_1:5:\the\ht0:\the\dp0}
\setbox0=\vbox{\hsize=1.040cm\parindent=0pt\linespread{1}\fontsize{10}{13.50}\selectfont\raggedright\noindent 中止\par}
\typeout{TCELL:p127_1:6:\the\ht0:\the\dp0}
\setbox0=\vbox{\hsize=8.840cm\parindent=0pt\linespread{1}\fontsize{10}{13.50}\selectfont\raggedright\noindent （1）结果未发生是中止行为的功劳。例如，甲用刀捅了乙，送乙到医院，救活。（2）结果未发生不是中止行为的功劳。例如，甲用刀捅了乙，叫来救护车，但乙已经被邻居送到医院，救活\par}
\typeout{TCELL:p127_1:7:\the\ht0:\the\dp0}
\setbox0=\vbox{\hsize=2.040cm\parindent=0pt\linespread{1}\fontsize{10}{13.50}\selectfont\raggedright\noindent $\times$\par}
\typeout{TCELL:p127_1:8:\the\ht0:\the\dp0}
\setbox0=\vbox{\hsize=3.040cm\parindent=0pt\linespread{1}\fontsize{10}{13.50}\selectfont\raggedright\noindent $\checkmark$\par}
\typeout{TCELL:p127_1:9:\the\ht0:\the\dp0}
\setbox0=\vbox{\hsize=1.040cm\parindent=0pt\linespread{1}\fontsize{10}{13.50}\selectfont\raggedright\noindent 未遂\par}
\typeout{TCELL:p127_1:10:\the\ht0:\the\dp0}
\setbox0=\vbox{\hsize=8.840cm\parindent=0pt\linespread{1}\fontsize{10}{13.50}\selectfont\raggedright\noindent 甲用刀捅了乙，扔包卫生纸后离去。乙被邻居救活\par}
\typeout{TCELL:p127_1:11:\the\ht0:\the\dp0}
\setbox0=\vbox{\hsize=2.040cm\parindent=0pt\linespread{1}\fontsize{10}{13.50}\selectfont\bfseries\centering\noindent 可能有效性\par}
\typeout{TCELL:p128_1:0:\the\ht0:\the\dp0}
\setbox0=\vbox{\hsize=3.040cm\parindent=0pt\linespread{1}\fontsize{10}{13.50}\selectfont\bfseries\centering\noindent 实际有效性\par}
\typeout{TCELL:p128_1:1:\the\ht0:\the\dp0}
\setbox0=\vbox{\hsize=1.040cm\parindent=0pt\linespread{1}\fontsize{10}{13.50}\selectfont\bfseries\centering\noindent 结论\par}
\typeout{TCELL:p128_1:2:\the\ht0:\the\dp0}
\setbox0=\vbox{\hsize=8.840cm\parindent=0pt\linespread{1}\fontsize{10}{13.50}\selectfont\bfseries\centering\noindent 范例\par}
\typeout{TCELL:p128_1:3:\the\ht0:\the\dp0}
\setbox0=\vbox{\hsize=2.040cm\parindent=0pt\linespread{1}\fontsize{10}{13.50}\selectfont\raggedright\noindent $\checkmark$\par}
\typeout{TCELL:p128_1:4:\the\ht0:\the\dp0}
\setbox0=\vbox{\hsize=3.040cm\parindent=0pt\linespread{1}\fontsize{10}{13.50}\selectfont\raggedright\noindent $\times$\par}
\typeout{TCELL:p128_1:5:\the\ht0:\the\dp0}
\setbox0=\vbox{\hsize=1.040cm\parindent=0pt\linespread{1}\fontsize{10}{13.50}\selectfont\raggedright\noindent 既遂\par}
\typeout{TCELL:p128_1:6:\the\ht0:\the\dp0}
\setbox0=\vbox{\hsize=8.840cm\parindent=0pt\linespread{1}\fontsize{10}{13.50}\selectfont\raggedright\noindent 甲用刀捅了乙，送乙到医院，抢救无效死亡\par}
\typeout{TCELL:p128_1:7:\the\ht0:\the\dp0}
\setbox0=\vbox{\hsize=2.040cm\parindent=0pt\linespread{1}\fontsize{10}{13.50}\selectfont\raggedright\noindent $\checkmark$\par}
\typeout{TCELL:p128_1:8:\the\ht0:\the\dp0}
\setbox0=\vbox{\hsize=3.040cm\parindent=0pt\linespread{1}\fontsize{10}{13.50}\selectfont\raggedright\noindent （—介入因素—）$\times$\par}
\typeout{TCELL:p128_1:9:\the\ht0:\the\dp0}
\setbox0=\vbox{\hsize=1.040cm\parindent=0pt\linespread{1}\fontsize{10}{13.50}\selectfont\raggedright\noindent 中止\par}
\typeout{TCELL:p128_1:10:\the\ht0:\the\dp0}
\setbox0=\vbox{\hsize=8.840cm\parindent=0pt\linespread{1}\fontsize{10}{13.50}\selectfont\raggedright\noindent 甲用刀捅了乙，送乙去医院，遇到车祸，车祸导致乙死亡\par}
\typeout{TCELL:p128_1:11:\the\ht0:\the\dp0}
\setbox0=\vbox{\hsize=5.240cm\parindent=0pt\linespread{1}\fontsize{10}{13.50}\selectfont\bfseries\centering\noindent 阶段与角色\par}
\typeout{TCELL:p136_1:0:\the\ht0:\the\dp0}
\setbox0=\vbox{\hsize=1.540cm\parindent=0pt\linespread{1}\fontsize{10}{13.50}\selectfont\bfseries\centering\noindent 实行者\par}
\typeout{TCELL:p136_1:1:\the\ht0:\the\dp0}
\setbox0=\vbox{\hsize=5.040cm\parindent=0pt\linespread{1}\fontsize{10}{13.50}\selectfont\bfseries\centering\noindent 教唆者、帮助者、另一实行者\par}
\typeout{TCELL:p136_1:2:\the\ht0:\the\dp0}
\setbox0=\vbox{\hsize=3.140cm\parindent=0pt\linespread{1}\fontsize{10}{13.50}\selectfont\bfseries\centering\noindent 结论\par}
\typeout{TCELL:p136_1:3:\the\ht0:\the\dp0}
\setbox0=\vbox{\hsize=5.240cm\parindent=0pt\linespread{1}\fontsize{10}{13.50}\selectfont\raggedright\noindent 制造违法事实阶段\par}
\typeout{TCELL:p136_1:4:\the\ht0:\the\dp0}
\setbox0=\vbox{\hsize=1.540cm\parindent=0pt\linespread{1}\fontsize{10}{13.50}\selectfont\raggedright\noindent $\checkmark$\par}
\typeout{TCELL:p136_1:5:\the\ht0:\the\dp0}
\setbox0=\vbox{\hsize=5.040cm\parindent=0pt\linespread{1}\fontsize{10}{13.50}\selectfont\raggedright\noindent $\checkmark$\par}
\typeout{TCELL:p136_1:6:\the\ht0:\the\dp0}
\setbox0=\vbox{\hsize=3.140cm\parindent=0pt\linespread{1}\fontsize{10}{13.50}\selectfont\raggedright\noindent 违法是一起的\par}
\typeout{TCELL:p136_1:7:\the\ht0:\the\dp0}
\setbox0=\vbox{\hsize=5.240cm\parindent=0pt\linespread{1}\fontsize{10}{13.50}\selectfont\raggedright\noindent 排除责任（谴责）阶段\par}
\typeout{TCELL:p136_1:8:\the\ht0:\the\dp0}
\setbox0=\vbox{\hsize=1.540cm\parindent=0pt\linespread{1}\fontsize{10}{13.50}\selectfont\raggedright\noindent $\times$\par}
\typeout{TCELL:p136_1:9:\the\ht0:\the\dp0}
\setbox0=\vbox{\hsize=5.040cm\parindent=0pt\linespread{1}\fontsize{10}{13.50}\selectfont\raggedright\noindent $\checkmark$\par}
\typeout{TCELL:p136_1:10:\the\ht0:\the\dp0}
\setbox0=\vbox{\hsize=3.140cm\parindent=0pt\linespread{1}\fontsize{10}{13.50}\selectfont\raggedright\noindent 责任是独立的\par}
\typeout{TCELL:p136_1:11:\the\ht0:\the\dp0}
\setbox0=\vbox{\hsize=2.640cm\parindent=0pt\linespread{1}\fontsize{10}{13.50}\selectfont\bfseries\centering\noindent 共同正犯的\\成立条件\par}
\typeout{TCELL:p137_1:0:\the\ht0:\the\dp0}
\setbox0=\vbox{\hsize=3.040cm\parindent=0pt\linespread{1}\fontsize{10}{13.50}\selectfont\centering\noindent 客观要件\par}
\typeout{TCELL:p137_1:1:\the\ht0:\the\dp0}
\setbox0=\vbox{\hsize=9.540cm\parindent=0pt\linespread{1}\fontsize{10}{13.50}\selectfont\centering\noindent 一起实施实行行为（意思联络+实行行为）\par}
\typeout{TCELL:p137_1:2:\the\ht0:\the\dp0}
\setbox0=\vbox{\hsize=3.040cm\parindent=0pt\linespread{1}\fontsize{10}{13.50}\selectfont\centering\noindent 主观要件\par}
\typeout{TCELL:p137_1:3:\the\ht0:\the\dp0}
\setbox0=\vbox{\hsize=9.540cm\parindent=0pt\linespread{1}\fontsize{10}{13.50}\selectfont\centering\noindent 故意（多数说）\par}
\typeout{TCELL:p137_1:4:\the\ht0:\the\dp0}
\setbox0=\vbox{\hsize=5.240cm\parindent=0pt\linespread{1}\fontsize{10}{13.50}\selectfont\bfseries\centering\noindent 案情\par}
\typeout{TCELL:p140_1:0:\the\ht0:\the\dp0}
\setbox0=\vbox{\hsize=1.740cm\parindent=0pt\linespread{1}\fontsize{10}{13.50}\selectfont\bfseries\centering\noindent 查证情况\par}
\typeout{TCELL:p140_1:1:\the\ht0:\the\dp0}
\setbox0=\vbox{\hsize=8.240cm\parindent=0pt\linespread{1}\fontsize{10}{13.50}\selectfont\bfseries\centering\noindent 结论\par}
\typeout{TCELL:p140_1:2:\the\ht0:\the\dp0}
\setbox0=\vbox{\hsize=5.240cm\parindent=0pt\linespread{1}\fontsize{10}{13.50}\selectfont\raggedright\noindent （1）故意的同时犯。甲、乙没有意思联络，故意杀人，同时向丙开枪，丙死亡\par}
\typeout{TCELL:p140_1:3:\the\ht0:\the\dp0}
\setbox0=\vbox{\hsize=5.240cm\parindent=0pt\linespread{1}\fontsize{10}{13.50}\selectfont\raggedright\noindent （2）过失的同时犯。甲、乙没有意思联络，各自过失导致子弹飞向丙，丙死亡\par}
\typeout{TCELL:p140_1:4:\the\ht0:\the\dp0}
\setbox0=\vbox{\hsize=5.240cm\parindent=0pt\linespread{1}\fontsize{10}{13.50}\selectfont\raggedright\noindent （3）故意的共同正犯。甲、乙有意思联络，共同故意，同时向丙开枪，丙死亡\par}
\typeout{TCELL:p140_1:5:\the\ht0:\the\dp0}
\setbox0=\vbox{\hsize=5.240cm\parindent=0pt\linespread{1}\fontsize{10}{13.50}\selectfont\raggedright\noindent （4）共同过失案：甲、乙有意思联络，共同过失，同时向丙开枪，丙死亡（打野猪案）\par}
\typeout{TCELL:p140_1:6:\the\ht0:\the\dp0}
\setbox0=\vbox{\hsize=1.740cm\parindent=0pt\linespread{1}\fontsize{10}{13.50}\selectfont\raggedright\noindent 事后查明，只有一枪致命，但无法查明致命一枪是谁打的\par}
\typeout{TCELL:p140_1:7:\the\ht0:\the\dp0}
\setbox0=\vbox{\hsize=8.240cm\parindent=0pt\linespread{1}\fontsize{10}{13.50}\selectfont\raggedright\noindent 甲、乙不构成共同正犯，需单独处理。根据存疑时有利于被告原则，甲、乙都构成故意杀人罪未遂，而非既遂\par}
\typeout{TCELL:p140_1:8:\the\ht0:\the\dp0}
\setbox0=\vbox{\hsize=8.240cm\parindent=0pt\linespread{1}\fontsize{10}{13.50}\selectfont\raggedright\noindent 甲、乙不构成共同正犯，需单独处理。根据存疑时有利于被告原则，甲、乙都不构成过失致人死亡罪，都无罪\par}
\typeout{TCELL:p140_1:9:\the\ht0:\the\dp0}
\setbox0=\vbox{\hsize=8.240cm\parindent=0pt\linespread{1}\fontsize{10}{13.50}\selectfont\raggedright\noindent 由于甲、乙是共同正犯，根据“部分实行、全部负责”原则，甲、乙对死亡结果都要负责，都构成故意杀人罪既遂\par}
\typeout{TCELL:p140_1:10:\the\ht0:\the\dp0}
\setbox0=\vbox{\hsize=8.240cm\parindent=0pt\linespread{1}\fontsize{10}{13.50}\selectfont\raggedright\noindent （1）根据要求故意且无罪说（少数说），甲、乙不构成共同正犯，基于“存疑时有利于被告”原则，都不构成过失致人死亡罪，都无罪。（2）根据要求故意且有罪说（多数说），甲、乙相互给对方的致人死亡的结果都有心理性贡献，因此对对方的致人死亡应负责。甲乙不构成共同正犯，各自构成过失致人死亡罪\par}
\typeout{TCELL:p140_1:11:\the\ht0:\the\dp0}
\setbox0=\vbox{\hsize=2.640cm\parindent=0pt\linespread{1}\fontsize{10}{13.50}\selectfont\bfseries\centering\noindent 间接正犯的\\成立条件\par}
\typeout{TCELL:p142_1:0:\the\ht0:\the\dp0}
\setbox0=\vbox{\hsize=3.040cm\parindent=0pt\linespread{1}\fontsize{10}{13.50}\selectfont\centering\noindent 客观要件\par}
\typeout{TCELL:p142_1:1:\the\ht0:\the\dp0}
\setbox0=\vbox{\hsize=9.540cm\parindent=0pt\linespread{1}\fontsize{10}{13.50}\selectfont\centering\noindent 引起+支配\par}
\typeout{TCELL:p142_1:2:\the\ht0:\the\dp0}
\setbox0=\vbox{\hsize=3.040cm\parindent=0pt\linespread{1}\fontsize{10}{13.50}\selectfont\centering\noindent 主观要件\par}
\typeout{TCELL:p142_1:3:\the\ht0:\the\dp0}
\setbox0=\vbox{\hsize=9.540cm\parindent=0pt\linespread{1}\fontsize{10}{13.50}\selectfont\centering\noindent 故意（故意支配他人）\par}
\typeout{TCELL:p142_1:4:\the\ht0:\the\dp0}
\setbox0=\vbox{\hsize=1.940cm\parindent=0pt\linespread{1}\fontsize{10}{13.50}\selectfont\bfseries\centering\noindent 案情\par}
\typeout{TCELL:p145_1:0:\the\ht0:\the\dp0}
\setbox0=\vbox{\hsize=4.140cm\parindent=0pt\linespread{1}\fontsize{10}{13.50}\selectfont\bfseries\centering\noindent 正犯的表现\par}
\typeout{TCELL:p145_1:1:\the\ht0:\the\dp0}
\setbox0=\vbox{\hsize=9.140cm\parindent=0pt\linespread{1}\fontsize{10}{13.50}\selectfont\bfseries\centering\noindent 共犯的结论\par}
\typeout{TCELL:p145_1:2:\the\ht0:\the\dp0}
\setbox0=\vbox{\hsize=1.940cm\parindent=0pt\linespread{1}\fontsize{10}{13.50}\selectfont\centering\noindent 甲教唆或帮助乙犯罪\par}
\typeout{TCELL:p145_1:3:\the\ht0:\the\dp0}
\setbox0=\vbox{\hsize=4.140cm\parindent=0pt\linespread{1}\fontsize{10}{13.50}\selectfont\raggedright\noindent 乙根本未去犯罪，无罪\par}
\typeout{TCELL:p145_1:4:\the\ht0:\the\dp0}
\setbox0=\vbox{\hsize=9.140cm\parindent=0pt\linespread{1}\fontsize{10}{13.50}\selectfont\centering\noindent 甲无罪\par}
\typeout{TCELL:p145_1:5:\the\ht0:\the\dp0}
\setbox0=\vbox{\hsize=4.140cm\parindent=0pt\linespread{1}\fontsize{10}{13.50}\selectfont\raggedright\noindent 乙在预备阶段，犯罪中止\par}
\typeout{TCELL:p145_1:6:\the\ht0:\the\dp0}
\setbox0=\vbox{\hsize=9.140cm\parindent=0pt\linespread{1}\fontsize{10}{13.50}\selectfont\raggedright\noindent （1）甲构成犯罪预备。（2）甲若有中止行为且消除贡献，则构成犯罪中止\par}
\typeout{TCELL:p145_1:7:\the\ht0:\the\dp0}
\setbox0=\vbox{\hsize=4.140cm\parindent=0pt\linespread{1}\fontsize{10}{13.50}\selectfont\raggedright\noindent 乙在预备阶段，犯罪预备\par}
\typeout{TCELL:p145_1:8:\the\ht0:\the\dp0}
\setbox0=\vbox{\hsize=9.140cm\parindent=0pt\linespread{1}\fontsize{10}{13.50}\selectfont\centering\noindent 甲构成犯罪预备\par}
\typeout{TCELL:p145_1:9:\the\ht0:\the\dp0}
\setbox0=\vbox{\hsize=4.140cm\parindent=0pt\linespread{1}\fontsize{10}{13.50}\selectfont\raggedright\noindent 乙在实行阶段，犯罪中止\par}
\typeout{TCELL:p145_1:10:\the\ht0:\the\dp0}
\setbox0=\vbox{\hsize=9.140cm\parindent=0pt\linespread{1}\fontsize{10}{13.50}\selectfont\raggedright\noindent （1）甲构成犯罪未遂。（2）甲若有中止行为且消除贡献，则构成犯罪中止\par}
\typeout{TCELL:p145_1:11:\the\ht0:\the\dp0}
\setbox0=\vbox{\hsize=4.140cm\parindent=0pt\linespread{1}\fontsize{10}{13.50}\selectfont\raggedright\noindent 乙在实行阶段，犯罪未遂\par}
\typeout{TCELL:p145_1:12:\the\ht0:\the\dp0}
\setbox0=\vbox{\hsize=9.140cm\parindent=0pt\linespread{1}\fontsize{10}{13.50}\selectfont\centering\noindent 甲构成犯罪未遂\par}
\typeout{TCELL:p145_1:13:\the\ht0:\the\dp0}
\setbox0=\vbox{\hsize=4.140cm\parindent=0pt\linespread{1}\fontsize{10}{13.50}\selectfont\raggedright\noindent 乙在实行阶段，犯罪既遂\par}
\typeout{TCELL:p145_1:14:\the\ht0:\the\dp0}
\setbox0=\vbox{\hsize=9.140cm\parindent=0pt\linespread{1}\fontsize{10}{13.50}\selectfont\raggedright\noindent （1）甲对既遂结果若有贡献，则构成犯罪既遂。（2）甲若有中止行为且消除贡献，则构成犯罪中止，对乙的既遂结果不用负责\par}
\typeout{TCELL:p145_1:15:\the\ht0:\the\dp0}
\setbox0=\vbox{\hsize=1.840cm\parindent=0pt\linespread{1}\fontsize{10}{13.50}\selectfont\bfseries\centering\noindent 帮助犯的\\成立条件\par}
\typeout{TCELL:p149_1:0:\the\ht0:\the\dp0}
\setbox0=\vbox{\hsize=2.140cm\parindent=0pt\linespread{1}\fontsize{10}{13.50}\selectfont\centering\noindent 客观要件\par}
\typeout{TCELL:p149_1:1:\the\ht0:\the\dp0}
\setbox0=\vbox{\hsize=3.940cm\parindent=0pt\linespread{1}\fontsize{10}{13.50}\selectfont\centering\noindent 促进正犯制造违法事实\par}
\typeout{TCELL:p149_1:2:\the\ht0:\the\dp0}
\setbox0=\vbox{\hsize=7.040cm\parindent=0pt\linespread{1}\fontsize{10}{13.50}\selectfont\raggedright\noindent 促进正犯制造违法行为，是帮助犯的成立条件\par}
\typeout{TCELL:p149_1:3:\the\ht0:\the\dp0}
\setbox0=\vbox{\hsize=7.040cm\parindent=0pt\linespread{1}\fontsize{10}{13.50}\selectfont\raggedright\noindent 促进正犯制造违法结果，是帮助犯的既遂条件\par}
\typeout{TCELL:p149_1:4:\the\ht0:\the\dp0}
\setbox0=\vbox{\hsize=2.140cm\parindent=0pt\linespread{1}\fontsize{10}{13.50}\selectfont\centering\noindent 主观要件\par}
\typeout{TCELL:p149_1:5:\the\ht0:\the\dp0}
\setbox0=\vbox{\hsize=11.240cm\parindent=0pt\linespread{1}\fontsize{10}{13.50}\selectfont\centering\noindent 故意促进\par}
\typeout{TCELL:p149_1:6:\the\ht0:\the\dp0}
\setbox0=\vbox{\hsize=1.340cm\parindent=0pt\linespread{1}\fontsize{10}{13.50}\selectfont\bfseries\centering\noindent 分类\par}
\typeout{TCELL:p151_2:0:\the\ht0:\the\dp0}
\setbox0=\vbox{\hsize=5.940cm\parindent=0pt\linespread{1}\fontsize{10}{13.50}\selectfont\bfseries\centering\noindent 案件事实\par}
\typeout{TCELL:p151_2:1:\the\ht0:\the\dp0}
\setbox0=\vbox{\hsize=1.940cm\parindent=0pt\linespread{1}\fontsize{10}{13.50}\selectfont\bfseries\centering\noindent 结论\par}
\typeout{TCELL:p151_2:2:\the\ht0:\the\dp0}
\setbox0=\vbox{\hsize=5.740cm\parindent=0pt\linespread{1}\fontsize{10}{13.50}\selectfont\bfseries\centering\noindent 称谓\par}
\typeout{TCELL:p151_2:3:\the\ht0:\the\dp0}
\setbox0=\vbox{\hsize=1.340cm\parindent=0pt\linespread{1}\fontsize{10}{13.50}\selectfont\centering\noindent 帮助\par}
\typeout{TCELL:p151_2:4:\the\ht0:\the\dp0}
\setbox0=\vbox{\hsize=5.940cm\parindent=0pt\linespread{1}\fontsize{10}{13.50}\selectfont\raggedright\noindent 提供不可能有用的帮助行为\par}
\typeout{TCELL:p151_2:5:\the\ht0:\the\dp0}
\setbox0=\vbox{\hsize=1.940cm\parindent=0pt\linespread{1}\fontsize{10}{13.50}\selectfont\centering\noindent 帮助者无罪\par}
\typeout{TCELL:p151_2:6:\the\ht0:\the\dp0}
\setbox0=\vbox{\hsize=5.740cm\parindent=0pt\linespread{1}\fontsize{10}{13.50}\selectfont\raggedright\noindent 对该案件事实，可称为“帮助未遂”，但不能称为“帮助犯未遂”“未遂的帮助犯”\par}
\typeout{TCELL:p151_2:7:\the\ht0:\the\dp0}
\setbox0=\vbox{\hsize=5.940cm\parindent=0pt\linespread{1}\fontsize{10}{13.50}\selectfont\raggedright\noindent 提供可能有用的帮助行为。包括两种情形：（1）促进正犯的实行行为，但没有促进正犯的既遂结果。（甲拿着乙提供有用的钥匙盗窃，操作不当打不开门，用其他办法盗窃既遂）（2）促进正犯的实行行为，但正犯没有既遂（甲用乙的钥匙打开了门，盗窃时被抓）\par}
\typeout{TCELL:p151_2:8:\the\ht0:\the\dp0}
\setbox0=\vbox{\hsize=1.940cm\parindent=0pt\linespread{1}\fontsize{10}{13.50}\selectfont\raggedright\noindent 帮助者成立帮助犯，但构成未遂\par}
\typeout{TCELL:p151_2:9:\the\ht0:\the\dp0}
\setbox0=\vbox{\hsize=5.740cm\parindent=0pt\linespread{1}\fontsize{10}{13.50}\selectfont\raggedright\noindent 第一，若侧重于表达帮助犯的犯罪形态，对该案件事实可称为“帮助犯未遂”。第二，若侧重于表达帮助犯帮正犯帮到哪个阶段和程度，对该案件事实可称为“未遂的帮助犯”，意指没有帮到正犯既遂的程度，只帮到正犯的实行阶段（也称为未遂阶段）\par}
\typeout{TCELL:p151_2:10:\the\ht0:\the\dp0}
\setbox0=\vbox{\hsize=1.340cm\parindent=0pt\linespread{1}\fontsize{10}{13.50}\selectfont\centering\noindent 教唆\par}
\typeout{TCELL:p151_2:11:\the\ht0:\the\dp0}
\setbox0=\vbox{\hsize=5.940cm\parindent=0pt\linespread{1}\fontsize{10}{13.50}\selectfont\raggedright\noindent 教唆实施不可能有危险的行为\par}
\typeout{TCELL:p151_2:12:\the\ht0:\the\dp0}
\setbox0=\vbox{\hsize=1.940cm\parindent=0pt\linespread{1}\fontsize{10}{13.50}\selectfont\centering\noindent 教唆者无罪\par}
\typeout{TCELL:p151_2:13:\the\ht0:\the\dp0}
\setbox0=\vbox{\hsize=5.740cm\parindent=0pt\linespread{1}\fontsize{10}{13.50}\selectfont\raggedright\noindent 对该案件事实，可称为“未遂的教唆”，但不能称为“未遂的教唆犯”\par}
\typeout{TCELL:p151_2:14:\the\ht0:\the\dp0}
\setbox0=\vbox{\hsize=5.940cm\parindent=0pt\linespread{1}\fontsize{10}{13.50}\selectfont\raggedright\noindent 教唆实施可能有危险的行为。正犯构成未遂\par}
\typeout{TCELL:p151_2:15:\the\ht0:\the\dp0}
\setbox0=\vbox{\hsize=1.940cm\parindent=0pt\linespread{1}\fontsize{10}{13.50}\selectfont\raggedright\noindent 教唆者成立教唆犯，但构成未遂\par}
\typeout{TCELL:p151_2:16:\the\ht0:\the\dp0}
\setbox0=\vbox{\hsize=5.740cm\parindent=0pt\linespread{1}\fontsize{10}{13.50}\selectfont\raggedright\noindent 对该案件事实，可称为“教唆犯未遂”，或“未遂的教唆犯”\par}
\typeout{TCELL:p151_2:17:\the\ht0:\the\dp0}
\setbox0=\vbox{\hsize=2.940cm\parindent=0pt\linespread{1}\fontsize{10}{13.50}\selectfont\bfseries\centering\noindent 前提事实\par}
\typeout{TCELL:p152_1:0:\the\ht0:\the\dp0}
\setbox0=\vbox{\hsize=4.940cm\parindent=0pt\linespread{1}\fontsize{10}{13.50}\selectfont\bfseries\centering\noindent 具体情形\par}
\typeout{TCELL:p152_1:1:\the\ht0:\the\dp0}
\setbox0=\vbox{\hsize=7.340cm\parindent=0pt\linespread{1}\fontsize{10}{13.50}\selectfont\bfseries\centering\noindent 结论\par}
\typeout{TCELL:p152_1:2:\the\ht0:\the\dp0}
\setbox0=\vbox{\hsize=2.940cm\parindent=0pt\linespread{1}\fontsize{10}{13.50}\selectfont\raggedright\noindent 甲准备拿着自己配的钥匙去盗窃丙家，乙知道后，送甲一把有用的钥匙，甲有了“双保险”的心理\par}
\typeout{TCELL:p152_1:3:\the\ht0:\the\dp0}
\setbox0=\vbox{\hsize=4.940cm\parindent=0pt\linespread{1}\fontsize{10}{13.50}\selectfont\raggedright\noindent 备而不用。甲先用自己的钥匙，打开了门，窃得财物\par}
\typeout{TCELL:p152_1:4:\the\ht0:\the\dp0}
\setbox0=\vbox{\hsize=7.340cm\parindent=0pt\linespread{1}\fontsize{10}{13.50}\selectfont\raggedright\noindent 乙的物理性帮助虽然没有发挥实际的帮助作用，但是起到心理性帮助效果，甲的心里有“双保险”的心理优势，故乙构成帮助犯既遂\par}
\typeout{TCELL:p152_1:5:\the\ht0:\the\dp0}
\setbox0=\vbox{\hsize=4.940cm\parindent=0pt\linespread{1}\fontsize{10}{13.50}\selectfont\raggedright\noindent 弃而不用。甲先用乙的钥匙，一时打不开门，便用自己的钥匙，打开了门，窃得财物\par}
\typeout{TCELL:p152_1:6:\the\ht0:\the\dp0}
\setbox0=\vbox{\hsize=7.340cm\parindent=0pt\linespread{1}\fontsize{10}{13.50}\selectfont\raggedright\noindent 乙的钥匙没有发挥实际的物理性帮助作用，也没有起到心理性的帮助效果，故乙不构成帮助犯既遂，而是未遂\par}
\typeout{TCELL:p152_1:7:\the\ht0:\the\dp0}
\setbox0=\vbox{\hsize=4.940cm\parindent=0pt\linespread{1}\fontsize{10}{13.50}\selectfont\raggedright\noindent 有用。甲先用乙的钥匙，打开了门，窃得财物\par}
\typeout{TCELL:p152_1:8:\the\ht0:\the\dp0}
\setbox0=\vbox{\hsize=7.340cm\parindent=0pt\linespread{1}\fontsize{10}{13.50}\selectfont\raggedright\noindent 乙的钥匙发挥了实际的物理性帮助，乙构成帮助犯既遂\par}
\typeout{TCELL:p152_1:9:\the\ht0:\the\dp0}
\setbox0=\vbox{\hsize=1.540cm\parindent=0pt\linespread{1}\fontsize{10}{13.50}\selectfont\bfseries\centering\noindent 角色\par}
\typeout{TCELL:p153_1:0:\the\ht0:\the\dp0}
\setbox0=\vbox{\hsize=3.040cm\parindent=0pt\linespread{1}\fontsize{10}{13.50}\selectfont\bfseries\centering\noindent 主观要件\par}
\typeout{TCELL:p153_1:1:\the\ht0:\the\dp0}
\setbox0=\vbox{\hsize=10.640cm\parindent=0pt\linespread{1}\fontsize{10}{13.50}\selectfont\bfseries\centering\noindent 范例\par}
\typeout{TCELL:p153_1:2:\the\ht0:\the\dp0}
\setbox0=\vbox{\hsize=3.040cm\parindent=0pt\linespread{1}\fontsize{10}{13.50}\selectfont\centering\noindent 都要求故意\par}
\typeout{TCELL:p153_1:3:\the\ht0:\the\dp0}
\setbox0=\vbox{\hsize=1.540cm\parindent=0pt\linespread{1}\fontsize{10}{13.50}\selectfont\centering\noindent 帮助犯\par}
\typeout{TCELL:p153_1:4:\the\ht0:\the\dp0}
\setbox0=\vbox{\hsize=10.640cm\parindent=0pt\linespread{1}\fontsize{10}{13.50}\selectfont\raggedright\noindent 甲明知乙醉酒并要开车，仍将车借给乙。乙醉驾不慎轧死丙。乙构成交通肇事罪。甲存在监管过失，构成交通肇事罪的单独犯罪。在此，不能说“甲构成交通肇事罪的帮助犯”\par}
\typeout{TCELL:p153_1:5:\the\ht0:\the\dp0}
\setbox0=\vbox{\hsize=1.540cm\parindent=0pt\linespread{1}\fontsize{10}{13.50}\selectfont\centering\noindent 教唆犯\par}
\typeout{TCELL:p153_1:6:\the\ht0:\the\dp0}
\setbox0=\vbox{\hsize=10.640cm\parindent=0pt\linespread{1}\fontsize{10}{13.50}\selectfont\raggedright\noindent 乙开车送领导甲去机场。甲为了赶时间，教唆乙闯红灯，认为没事。乙照办，不慎撞死人。乙构成交通肇事罪。甲存在指挥过失，构成交通肇事罪的单独犯罪。在此不能说“甲构成交通肇事罪的教唆犯”\par}
\typeout{TCELL:p153_1:7:\the\ht0:\the\dp0}
\setbox0=\vbox{\hsize=1.540cm\parindent=0pt\linespread{1}\fontsize{10}{13.50}\selectfont\centering\noindent 间接正犯\par}
\typeout{TCELL:p153_1:8:\the\ht0:\the\dp0}
\setbox0=\vbox{\hsize=10.640cm\parindent=0pt\linespread{1}\fontsize{10}{13.50}\selectfont\raggedright\noindent 医生欺骗护士给病人打了毒针，导致病人死亡。护士构成医疗事故罪，但不能说医生构成医疗事故罪的间接正犯。医生构成故意杀人罪的间接正犯\par}
\typeout{TCELL:p153_1:9:\the\ht0:\the\dp0}
\setbox0=\vbox{\hsize=1.540cm\parindent=0pt\linespread{1}\fontsize{10}{13.50}\selectfont\centering\noindent 共同正犯\par}
\typeout{TCELL:p153_1:10:\the\ht0:\the\dp0}
\setbox0=\vbox{\hsize=10.640cm\parindent=0pt\linespread{1}\fontsize{10}{13.50}\selectfont\raggedright\noindent 甲、乙有意思联络，误将丙当作野猪打死，无法查明谁打的致命一枪。（1）根据要求故意且无罪说（少数说），甲乙不构成共同正犯，甲乙都无罪。（2）根据要求故意且有罪说（多数说），甲乙不构成共同正犯，甲乙各自构成过失致人死亡罪\par}
\typeout{TCELL:p153_1:11:\the\ht0:\the\dp0}
\setbox0=\vbox{\hsize=2.040cm\parindent=0pt\linespread{1}\fontsize{10}{13.50}\selectfont\bfseries\centering\noindent 参与时间\par}
\typeout{TCELL:p156_1:0:\the\ht0:\the\dp0}
\setbox0=\vbox{\hsize=3.940cm\parindent=0pt\linespread{1}\fontsize{10}{13.50}\selectfont\bfseries\centering\noindent 结论\par}
\typeout{TCELL:p156_1:1:\the\ht0:\the\dp0}
\setbox0=\vbox{\hsize=9.240cm\parindent=0pt\linespread{1}\fontsize{10}{13.50}\selectfont\bfseries\centering\noindent 范例\par}
\typeout{TCELL:p156_1:2:\the\ht0:\the\dp0}
\setbox0=\vbox{\hsize=2.040cm\parindent=0pt\linespread{1}\fontsize{10}{13.50}\selectfont\centering\noindent 事前共谋，临时迟到\par}
\typeout{TCELL:p156_1:3:\the\ht0:\the\dp0}
\setbox0=\vbox{\hsize=3.940cm\parindent=0pt\linespread{1}\fontsize{10}{13.50}\selectfont\centering\noindent 成立共同犯罪\par}
\typeout{TCELL:p156_1:4:\the\ht0:\the\dp0}
\setbox0=\vbox{\hsize=9.240cm\parindent=0pt\linespread{1}\fontsize{10}{13.50}\selectfont\raggedright\noindent 甲乙共谋明晚抢劫丙，甲先到，重伤丙，乙赶到，拿走丙财物。甲乙构成抢劫罪（致人重伤）的共同犯罪（2018 年试题）\par}
\typeout{TCELL:p156_1:5:\the\ht0:\the\dp0}
\setbox0=\vbox{\hsize=2.040cm\parindent=0pt\linespread{1}\fontsize{10}{13.50}\selectfont\centering\noindent 中途参与\par}
\typeout{TCELL:p156_1:6:\the\ht0:\the\dp0}
\setbox0=\vbox{\hsize=3.940cm\parindent=0pt\linespread{1}\fontsize{10}{13.50}\selectfont\centering\noindent 成立承继的共同犯罪\par}
\typeout{TCELL:p156_1:7:\the\ht0:\the\dp0}
\setbox0=\vbox{\hsize=9.240cm\parindent=0pt\linespread{1}\fontsize{10}{13.50}\selectfont\raggedright\noindent 甲抢劫丙，重伤丙后，乙才参与，拿走丙的财物，交给甲。甲构成抢劫罪（致人重伤）。乙构成甲的抢劫罪的承继共同犯罪，但对重伤不负责，因为重伤是乙参与之前甲制造的（2007 年第 53 题）②\par}
\typeout{TCELL:p156_1:8:\the\ht0:\the\dp0}
\setbox0=\vbox{\hsize=2.040cm\parindent=0pt\linespread{1}\fontsize{10}{13.50}\selectfont\centering\noindent 事后参与（事后帮）③\par}
\typeout{TCELL:p156_1:9:\the\ht0:\the\dp0}
\setbox0=\vbox{\hsize=3.940cm\parindent=0pt\linespread{1}\fontsize{10}{13.50}\selectfont\raggedright\noindent 帮助对象是犯罪人，构成窝藏、包庇罪\par}
\typeout{TCELL:p156_1:10:\the\ht0:\the\dp0}
\setbox0=\vbox{\hsize=9.240cm\parindent=0pt\linespread{1}\fontsize{10}{13.50}\selectfont\raggedright\noindent 甲抢劫到丙的财物后，乙骑车路过，载甲离去。乙构成窝藏罪\par}
\typeout{TCELL:p156_1:11:\the\ht0:\the\dp0}
\setbox0=\vbox{\hsize=3.940cm\parindent=0pt\linespread{1}\fontsize{10}{13.50}\selectfont\raggedright\noindent 帮助对象是赃物，构成掩饰、隐瞒犯罪所得罪\par}
\typeout{TCELL:p156_1:12:\the\ht0:\the\dp0}
\setbox0=\vbox{\hsize=9.240cm\parindent=0pt\linespread{1}\fontsize{10}{13.50}\selectfont\raggedright\noindent 甲抢劫到丙的财物后，乙帮助销赃。乙构成掩饰、隐瞒犯罪所得罪\par}
\typeout{TCELL:p156_1:13:\the\ht0:\the\dp0}
\setbox0=\vbox{\hsize=3.940cm\parindent=0pt\linespread{1}\fontsize{10}{13.50}\selectfont\raggedright\noindent 帮助对象是证据，构成帮助毁灭、伪造证据罪\par}
\typeout{TCELL:p156_1:14:\the\ht0:\the\dp0}
\setbox0=\vbox{\hsize=9.240cm\parindent=0pt\linespread{1}\fontsize{10}{13.50}\selectfont\raggedright\noindent 甲抢劫到丙的财物后，乙帮助销毁抢劫用的工具。乙构成帮助毁灭证据罪\par}
\typeout{TCELL:p156_1:15:\the\ht0:\the\dp0}
\setbox0=\vbox{\hsize=2.540cm\parindent=0pt\linespread{1}\fontsize{10}{13.50}\selectfont\bfseries\centering\noindent 角色\par}
\typeout{TCELL:p157_1:0:\the\ht0:\the\dp0}
\setbox0=\vbox{\hsize=12.940cm\parindent=0pt\linespread{1}\fontsize{10}{13.50}\selectfont\bfseries\centering\noindent 要求\par}
\typeout{TCELL:p157_1:1:\the\ht0:\the\dp0}
\setbox0=\vbox{\hsize=2.540cm\parindent=0pt\linespread{1}\fontsize{10}{13.50}\selectfont\raggedright\noindent 教唆犯\par}
\typeout{TCELL:p157_1:2:\the\ht0:\the\dp0}
\setbox0=\vbox{\hsize=12.940cm\parindent=0pt\linespread{1}\fontsize{10}{13.50}\selectfont\raggedright\noindent 消除心理性贡献（打消犯意）\par}
\typeout{TCELL:p157_1:3:\the\ht0:\the\dp0}
\setbox0=\vbox{\hsize=2.540cm\parindent=0pt\linespread{1}\fontsize{10}{13.50}\selectfont\raggedright\noindent 帮助犯\par}
\typeout{TCELL:p157_1:4:\the\ht0:\the\dp0}
\setbox0=\vbox{\hsize=12.940cm\parindent=0pt\linespread{1}\fontsize{10}{13.50}\selectfont\raggedright\noindent 消除物理性、心理性的贡献\par}
\typeout{TCELL:p157_1:5:\the\ht0:\the\dp0}
\setbox0=\vbox{\hsize=2.540cm\parindent=0pt\linespread{1}\fontsize{10}{13.50}\selectfont\raggedright\noindent 共同正犯\par}
\typeout{TCELL:p157_1:6:\the\ht0:\the\dp0}
\setbox0=\vbox{\hsize=12.940cm\parindent=0pt\linespread{1}\fontsize{10}{13.50}\selectfont\raggedright\noindent 消除物理性、心理性的贡献。\\在预备阶段，要消除心理性贡献，需要明确告知退出意思。\\在实行阶段，要消除贡献，原则上要求有效防止既遂结果发生\par}
\typeout{TCELL:p157_1:7:\the\ht0:\the\dp0}
\setbox0=\vbox{\hsize=8.040cm\parindent=0pt\linespread{1}\fontsize{10}{13.50}\selectfont\bfseries\centering\noindent 结合类型\par}
\typeout{TCELL:p161_1:0:\the\ht0:\the\dp0}
\setbox0=\vbox{\hsize=7.440cm\parindent=0pt\linespread{1}\fontsize{10}{13.50}\selectfont\bfseries\centering\noindent 处理结论\par}
\typeout{TCELL:p161_1:1:\the\ht0:\the\dp0}
\setbox0=\vbox{\hsize=8.040cm\parindent=0pt\linespread{1}\fontsize{10}{13.50}\selectfont\raggedright\noindent 有身份者是身份犯的实行犯+无身份者是共犯\par}
\typeout{TCELL:p161_1:2:\the\ht0:\the\dp0}
\setbox0=\vbox{\hsize=7.440cm\parindent=0pt\linespread{1}\fontsize{10}{13.50}\selectfont\raggedright\noindent 按实行犯来定罪（也即按有身份者定罪）（例 1）\par}
\typeout{TCELL:p161_1:3:\the\ht0:\the\dp0}
\setbox0=\vbox{\hsize=8.040cm\parindent=0pt\linespread{1}\fontsize{10}{13.50}\selectfont\raggedright\noindent 无身份者是非身份犯的实行犯+有身份者是共犯\par}
\typeout{TCELL:p161_1:4:\the\ht0:\the\dp0}
\setbox0=\vbox{\hsize=7.440cm\parindent=0pt\linespread{1}\fontsize{10}{13.50}\selectfont\raggedright\noindent 按实行犯来定罪（也即按无身份者定罪）（例 2）\par}
\typeout{TCELL:p161_1:5:\the\ht0:\the\dp0}
\setbox0=\vbox{\hsize=8.040cm\parindent=0pt\linespread{1}\fontsize{10}{13.50}\selectfont\raggedright\noindent A 身份者是 A 身份犯的实行犯+B 身份者是 B 身份犯的实行犯\par}
\typeout{TCELL:p161_1:6:\the\ht0:\the\dp0}
\setbox0=\vbox{\hsize=7.440cm\parindent=0pt\linespread{1}\fontsize{10}{13.50}\selectfont\raggedright\noindent 按实行犯来定罪，存在想象竞合，择一重罪论处（例 3）\par}
\typeout{TCELL:p161_1:7:\the\ht0:\the\dp0}
\setbox0=\vbox{\hsize=1.940cm\parindent=0pt\linespread{1}\fontsize{10}{13.50}\selectfont\bfseries\centering\noindent 作为义务人（甲）\par}
\typeout{TCELL:p162_1:0:\the\ht0:\the\dp0}
\setbox0=\vbox{\hsize=2.140cm\parindent=0pt\linespread{1}\fontsize{10}{13.50}\selectfont\bfseries\centering\noindent 非作为义务人（乙）\par}
\typeout{TCELL:p162_1:1:\the\ht0:\the\dp0}
\setbox0=\vbox{\hsize=5.340cm\parindent=0pt\linespread{1}\fontsize{10}{13.50}\selectfont\bfseries\centering\noindent 范例\par}
\typeout{TCELL:p162_1:2:\the\ht0:\the\dp0}
\setbox0=\vbox{\hsize=5.540cm\parindent=0pt\linespread{1}\fontsize{10}{13.50}\selectfont\bfseries\centering\noindent 结论\par}
\typeout{TCELL:p162_1:3:\the\ht0:\the\dp0}
\setbox0=\vbox{\hsize=1.940cm\parindent=0pt\linespread{1}\fontsize{10}{13.50}\selectfont\raggedright\noindent 实行\par}
\typeout{TCELL:p162_1:4:\the\ht0:\the\dp0}
\setbox0=\vbox{\hsize=2.140cm\parindent=0pt\linespread{1}\fontsize{10}{13.50}\selectfont\raggedright\noindent 帮助、教唆\par}
\typeout{TCELL:p162_1:5:\the\ht0:\the\dp0}
\setbox0=\vbox{\hsize=5.340cm\parindent=0pt\linespread{1}\fontsize{10}{13.50}\selectfont\raggedright\noindent 1. 甲不想履行扶养义务，要将老父亲丙遗弃，让乙帮忙。乙照办。2. 甲不慎导致丙受伤，想救丙。乙教唆甲不要救。甲便不救\par}
\typeout{TCELL:p162_1:6:\the\ht0:\the\dp0}
\setbox0=\vbox{\hsize=5.540cm\parindent=0pt\linespread{1}\fontsize{10}{13.50}\selectfont\raggedright\noindent 甲构成不作为犯的实行犯，乙构成帮助犯、教唆犯（作为义务人是一种身份，乙没这个身份，故不可能构成不作为犯的实行犯）\par}
\typeout{TCELL:p162_1:7:\the\ht0:\the\dp0}
\setbox0=\vbox{\hsize=1.940cm\parindent=0pt\linespread{1}\fontsize{10}{13.50}\selectfont\bfseries\centering\noindent 作为义务人（甲）\par}
\typeout{TCELL:p163_1:0:\the\ht0:\the\dp0}
\setbox0=\vbox{\hsize=2.140cm\parindent=0pt\linespread{1}\fontsize{10}{13.50}\selectfont\bfseries\centering\noindent 非作为义务人（乙）\par}
\typeout{TCELL:p163_1:1:\the\ht0:\the\dp0}
\setbox0=\vbox{\hsize=5.340cm\parindent=0pt\linespread{1}\fontsize{10}{13.50}\selectfont\bfseries\centering\noindent 范例\par}
\typeout{TCELL:p163_1:2:\the\ht0:\the\dp0}
\setbox0=\vbox{\hsize=5.540cm\parindent=0pt\linespread{1}\fontsize{10}{13.50}\selectfont\bfseries\centering\noindent 结论\par}
\typeout{TCELL:p163_1:3:\the\ht0:\the\dp0}
\setbox0=\vbox{\hsize=1.940cm\parindent=0pt\linespread{1}\fontsize{10}{13.50}\selectfont\centering\noindent 帮助\par}
\typeout{TCELL:p163_1:4:\the\ht0:\the\dp0}
\setbox0=\vbox{\hsize=2.140cm\parindent=0pt\linespread{1}\fontsize{10}{13.50}\selectfont\centering\noindent 实行\par}
\typeout{TCELL:p163_1:5:\the\ht0:\the\dp0}
\setbox0=\vbox{\hsize=5.340cm\parindent=0pt\linespread{1}\fontsize{10}{13.50}\selectfont\raggedright\noindent 乙用欺诈合同欲诈骗丙公司，让丙公司的法务甲睁只眼、闭只眼。甲答应照办\par}
\typeout{TCELL:p163_1:6:\the\ht0:\the\dp0}
\setbox0=\vbox{\hsize=5.540cm\parindent=0pt\linespread{1}\fontsize{10}{13.50}\selectfont\raggedright\noindent 乙构成诈骗罪的实行犯，甲构成不作为的帮助犯\par}
\typeout{TCELL:p163_1:7:\the\ht0:\the\dp0}
\setbox0=\vbox{\hsize=4.340cm\parindent=0pt\linespread{1}\fontsize{10}{13.50}\selectfont\raggedright\noindent 作为义务的来源根据就是总论的“实质二分说”。常考的是先行行为。判断标准：先行行为若导致他人陷入危险状态，则负有保护义务（保护人身法益和财产法益的义务）\par}
\typeout{TCELL:p163_1:8:\the\ht0:\the\dp0}
\setbox0=\vbox{\hsize=5.340cm\parindent=0pt\linespread{1}\fontsize{10}{13.50}\selectfont\raggedright\noindent 甲乙欲实施故意伤害罪，共同将丙女打晕。乙又要强奸丙或拿财物。甲看到，不阻止\par}
\typeout{TCELL:p163_1:9:\the\ht0:\the\dp0}
\setbox0=\vbox{\hsize=5.540cm\parindent=0pt\linespread{1}\fontsize{10}{13.50}\selectfont\raggedright\noindent 甲构成强奸罪的不作为的帮助犯、盗窃罪的不作为的帮助犯\par}
\typeout{TCELL:p163_1:10:\the\ht0:\the\dp0}
\mbox{}
\end{document}
